% r03_fill_K3a: proof of (K_0), part IV: the comparison principle, the two-flow bound, the constant-background reduction.

\subsubsection*{The comparison principle and the two-flow bound}

A \emph{renormalised family} is a finite sum of King's renormalised graphs \cite[\S3.5]{K1}, Ba{\l}aban's classes of \cite[\S\S2--3]{Ba3} (\Cref{def:a07:class}), localised as in \cite[(1.33)]{Ba3}. It is in \emph{per-flow normal form} when, for each flow separately and before any element is mixed, every counterterm graph is expanded in the lines of the flow \cite[(3.76)]{K1}, every Taylor subtraction \cite[(3.83)]{K1} and Ward-identity rewriting has been carried out with that flow's exact identities, and the exact cancellations between members with identical element sets (\cite[(3.6)/(3.7), (2.21b)/(2.21g), (2.18)]{Ba3}) are co-switched, the identified elements of both members being switched together. Its \emph{elements} are lines, vertex functions and external fields.

\begin{proposition}[comparison principle]\label{rot:prop:cp}
Let $\cR$ be a renormalised family in per-flow normal form with $n_\cR$ elements and $E^F(\cR)$ its value along flow $F$. Then, with $\theta=\min(\frac12,\theta_1)$ and $\theta_1$ of \Cref{rot:thm:H1S},
\[
|E^X(\cR)-E^Y(\cR)|\le Cn_\cR\,p_r^2\delta_r^{-1/2}\,\eta'^{\theta}(1+k)^p\cdot[\text{right side of \cite[(1.33)]{Ba3} for }\cR],
\]
the prefactor $p_r^2\delta_r^{-1/2}$ being that of the contour elements at $e>0$ (absent at $e=0$).
\end{proposition}
\begin{proof}
By \Cref{rot:lem:vert}(i) both values are sums over $X$-vertices. Order the elements consistently across members and telescope: $E^X-E^Y=\sum_iE[e_1,\dots,e_{i-1}\text{ in $X$-version};\ e_i\to e_i^X-e_i^Y;\ e_{i+1},\dots\text{ in $Y$-version}]$. In the $i$-th term the element $e^X_i-e^Y_i$, divided by $C\eta'^\theta$, obeys the bounds of its type weakened by the degree $\theta$: fine lines by \Cref{rot:cor:unif}(a), coarse lines by (b) ($\eta^{1/2}\le\sqrt5\eta'^{1/2}$), vertex functions by \Cref{rot:lem:vert}(ii) with $\eta'=\eta'^\theta\eta'^{1-\theta}$, external fields $H^F\phi_r$ and their covariant derivatives by \Cref{rot:lem:H} and its derivative and H\"older versions; at $e>0$, contour lines \cite[(3.27)]{K1} by \Cref{rot:cor:unif}(c) (gain $(\eta'/\ell_j)^{1/2}$), double contour lines by \Cref{rot:prop:BD}(b) (the same relative gain, size $C\ell_j$), contour vertices \cite[(3.24)--(3.26)]{K1} by \Cref{rot:prop:BD}(b) (relative gain $\eta^{1/2}$, size $Ce_rp_r^2$) and the constant-background phases \cite[(3.46)]{K1} by \Cref{rot:prop:BD}(b) ($\le e\,p_r\lmu^{-1}c_\tau\delta_r^{-1/2}\eta'$), whence the prefactor. The weakened-element bound for each term is \Cref{rot:thm:H1S}, whose proof runs this telescoping with the subtraction rule described there.
\end{proof}

\begin{theorem}[the two-flow comparison bound]\label{rot:thm:H1S}
Let $\alpha\in(0,1)$ be the H\"older exponent of \cite[(3.8)]{K1} and put $\theta_1\defeq\alpha^2/(4(1+\alpha)^2)$. There are $p=p(\bar n)$ and $C$ such that for every renormalised family $\cR$ of order $\le\bar n$ in per-flow normal form (the $h$-families of \Cref{rot:lem:S1} other than $\cT_h$, and the families of step (5) of the proof of \Cref{rot:thm:A} at $e>0$ after the constant-background reduction of \Cref{rot:lem:CB}), for $0<\theta\le\theta_1$, uniformly in $r$, $k$, the torus and the masses, the two-flow difference obeys the bound of \Cref{rot:prop:cp}: in every term of its telescoping, the element replaced by the difference $e^X_i-e^Y_i$, which obeys the bounds of its type weakened by the degree $\theta$, \cite[(1.33)]{Ba3} holds with the extra factor $C(1+k)^p\eta'^{\theta}$ and the prefactor stated there. The exponent $\theta_1$ does not depend on $\bar n$, and any $\theta>0$ with any polynomial factor suffices for the consumers (\Cref{rot:lem:r}).
\end{theorem}
The proof is given at the end of this part, after the constant-background reduction (\Cref{rot:lem:CB}), which it uses at $e>0$. Its inputs are \Cref{rot:prop:BD}, \cite[Thm.~3.3]{K1} for the standard flow, \cite[Prop.~2.1, (2.8)]{Ba3} with the ordering device (2.7) there, the zero-momentum counterterm conditions \cite[(3.33)]{K1} through order four with the vanishing of the one-loop vacuum-polarisation moments \cite[(2.23)--(2.28)]{Ba3}, the local limit law for the simple random walk \cite[Thm.~2.5.6]{LL10}, and the $h$-census of \Cref{rot:lem:S1}. In outline: every line is written as one flow-independent lattice function sliced by the heat kernel; degrees are read per ordering of the scale-summed elements; each unit (a self-energy, vacuum-polarisation or triangle insertion) is bundled with one of its legs into a single element, and one subtraction rule is applied, in each flow's own frame, to every connected component of degree zero that forms in an initial segment, its local terms being common numbers of both flows; after the rule every component containing the weakened element has degree $\ge\min(\alpha,1)$, and the scale sums give the gain $\eta'^{\theta_1}$.

\begin{lemma}[the constant-background reduction]\label{rot:lem:CB}
In the localised form \cite[(3.55)]{K1} every family has its external vertices in a cube $D$ of side $\le Cp_r$. Let $D''$ be the union of the level-$r$ blocks within $p$ of $D$, $x_0\in D$, $B\defeq A_r(x_0)$ and $\tilde A\defeq A^{(r)}-B$ on all of $T$. For each flow $F$ the reduction runs on the infinite fine lattice $\Lambda_F$ of level $r$ with the $P_T$-periodic lifts $A_{\rm lift}$, $\tilde A_{\rm lift}$ (no periodic cell is formed); $F_T$ is the lifted fundamental cube centred at $D$, $F_1\defeq\{\dist(\cdot,D'')\le p/3\}$, and the vertex region $F_v$ is the union of the level-$r$ blocks within $p/3-3$ of $D''$. The sup-systole of $P_T$ in the unit scale of level $r$ is $s_T\ge40MJ\,5^{\lfloor\vr k\rfloor}\ge Cp_r$ for $k\ge k_0$.
\begin{enumerate}[label=\textup{(\arabic*)},leftmargin=*]
\item[(0)] \textup{(P)} For $x,y\in F_T$, $G_r(T,A^{(r)})(x,y)=\sum_{\gamma\in P_T}G_r(\Lambda_F,A_{\rm lift})(x,y+\gamma)$, and likewise for $D_AG_r$, $Q_r$ and $H^F$; the image part $\delta^{\rm im}$ (the terms $\gamma\ne0$) obeys $|(\delta^{\rm im}f)(x)|,|(D_A\delta^{\rm im}f)(x)|\le Ce^{-\delta_0s_T/4}\|f\|_\infty$ for $\supp f\subset B(x,s_T/4)$. \textup{(G)} On $\Lambda_F$ the constant $B$ is removed exactly by the gauge function $B\cdot x$, leaving King's phases \cite[(3.46)]{K1} at the vertices. \textup{(R)} $|\tilde A_r(w)|\le p_r(d_T(w,x_0)+1)$ and $|\partial\tilde A^{(r)}|\le Cp_r$ on all of $T$, so $\tilde A_{\rm lift}$, $\vartheta\tilde A_{\rm lift}$ and $B_{\rm bg}\defeq(1-\vartheta)\tilde A_{\rm lift}$ are regular. \textup{(N)} If $A_1$, $A_2$ are regular on $\Lambda_F$ and $A_1=A_2$ on the bonds inside $F_1$, then $\delta\defeq G_r(\Lambda_F,A_1)-G_r(\Lambda_F,A_2)$ obeys $|\delta(x,y)|\le Ce^{-\delta_1(|x-y|+\dist(x,F_1^c)+\dist(y,F_1^c))}$ with $C$ independent of $F_1$, and the same for $D_A\delta$ when the bond at $x$ lies inside $F_1$; in particular $G_r(\Lambda_F,B_{\rm bg})=G_r(\Lambda_F,0)+\delta^{\rm bg}$.
\item For each flow $F\in\{X,\tau Y\}$ and each renormalised family $\cR$ at background $A^{(r)}$ on $\supp\chi_r$: $\cR^F(A^{(r)})=\sum_{q\le q_0}\sum_{w_1,\dots,w_q}A_r(w_1)\cdots A_r(w_q)\,\cR^F_{q;w}(A=0)\prod_vU^F(B(\Gamma_v))+\cR^F_{\delta G}+\cR^F_{\rm rem}$, where $\cR^F_{q;w}$ is a renormalised family in per-flow normal form whose lines are the $A=0$ lines of the flow ($G^F_r(0)=C-\Gamma^F_r(0)$, $H^F$, $(A^F)^{-1}$) and whose vertices are \cite[(3.18)--(3.26)]{K1} at $B=0$ with the explicit factors $\eta$ of the pin and contour vertices; the coarse-field factors $A_r(w_i)$ are the same functions of the common coarse field for both flows; $U^F(B(\Gamma_v))=\exp[\mathrm ie(5^r\eps)^{1/2}qB(\Gamma^F_v)]$ are the constant-background phases \cite[(3.46)]{K1}; $\cR^F_{\delta G}$ collects the terms with a line $\delta^{\rm im}$ or $\delta^{\rm bg}$ or with a vertex outside $F_v$, and $\cR^F_{\rm rem}$ the Taylor remainders of \cite[(3.48), (3.52)]{K1}.
\item $|\cR^F_{\delta G}|+|\cR^F_{\rm rem}|\le C(\delta_r^\sigma+e^{-\delta'p(\delta_r)})|T|$ per flow.
\item The two-flow difference of the leading sum is controlled element by element: the $A=0$ lines by \Cref{rot:cor:unif} and the difference elements $\Delta\Gamma$, $\Delta H$, $\Delta(A^{-1})$ of size $O(\eta'^{1/2})$; for a constant field $B$ the line integral depends only on the endpoints, $B(\Gamma^F_v)=B\cdot(c(v)-v)$, so $U^X(B(\Gamma_v))$ and $U^{\tau Y}(B(\Gamma_{\tau v}))$ differ by the phase of $e(5^r\eps)^{1/2}B\cdot(\tau v-v)$, $|\tau v-v|\le c_\tau\eta'$, i.e.\ by $\le ep_r\lmu^{-1}c_\tau\delta_r^{-1/2}\eta'$ (\cite[(3.2)]{K1}); the contour vertices \cite[(3.24)--(3.26)]{K1} at $B=0$ differ entrywise by a relative $\eta^{1/2}$ and the contour lines, one-ended and double, obey the slice-$j$ bounds in both flows with mismatch $(\eta'/\ell_j)^{1/2}$ (\Cref{rot:prop:BD}(b)); every connected subgraph containing a contour line has Ba{\l}aban's degree $\ge1$; the explicit factors $A_r(w_i)$ are common.
\end{enumerate}
\end{lemma}
\begin{proof}
(1) is King's expansion \cite[(3.43)--(3.55)]{K1} applied to each flow separately; it uses only the gauge covariance of the flow's lines (the hybrid flow is gauge covariant: its block average is the frame-covariant photon average of \Cref{rot:rem:frameav} and \Cref{rot:lem:tile}, so a constant $B$ is mapped to the coarse constant $R^{-1}B$ for the $Y$ flow and $B$ for the $X$ flow, and the gauge transformation $A^{(r)}\mapsto A^{(r)}-B$ acts on both by the same unitary at each vertex), the universal-cover form (0) of the localisation, and the Taylor expansion around $\vartheta\tilde A_{\rm lift}=0$ with background $B_{\rm bg}$, \cite[(3.47)--(3.54)]{K1}, exact algebra in which the lines become $A=0$ lines and the powers of $\tilde A$ become explicit coarse-field factors sampled at block points. The paired contours (\Cref{rot:lem:pair}) make the contour phases of the two flows differ by the flux through an area $O(\eta')$ (\Cref{rot:lem:flux}), so that the contour vertices and lines are difference elements with the gains of (3) rather than $O(1)$ mismatches. (0): the block systems of both flows, tiles and fixed staircases included, are $P_T$-invariant (\Cref{rot:lem:adm}(iii) and the in-tile convention of \Cref{rot:def:flows}), so the operator $M_{A_{\rm lift}}+a_rQ_F^*Q_F$ on $\Lambda_F$ commutes with $P_T$; its inverse has a summable kernel (the bounds of (2)), its periodisation is a right inverse of the torus operator, which is positive definite, and so equals its inverse; for $|x-y|\le s_T/2$ and $\gamma\ne0$, $|x-y-\gamma|\ge|\gamma|/2$, whence the image bound. (G): on $\Lambda_F$ the gauge function is single-valued and both flows are gauge covariant for every gauge function. (R): \cite[(3.2)]{K1} on all of $T$, $\sum_z\nabla H^{\rm ph}(x,z)=0$ with the $r$-uniform decay of $\nabla H^{\rm ph}$, and the product rule with $|\nabla\vartheta|\le3/p$. (N): the random-walk representation of \Cref{rot:prop:BD}(a) with one domain: chains whose cells lie with a fixed margin inside $F_1$ are the same for $A_1$ and $A_2$, every other chain of length $n$ has $\sqrt3M_{\rm w}n\ge\dist(x,F_1^c)+\dist(\supp f,F_1^c)-c$ and $\sqrt3M_{\rm w}n\ge|x-\supp f|-c$, and the average of the two gives the rate $\delta_1=(2\sqrt3M_{\rm w})^{-1}$. Vertices in $F_v$ see $B_{\rm bg}=0$ on every bond and in-block contour they touch, so King's ``taking $A'=\theta\tilde A$ and $B=0$'' \cite[p.~662]{K1} holds there literally. (2) is King's bound \cite[p.~665]{K1} with the kernel bounds at the regular fields of (R): a $\delta^{\rm bg}$ line with an end within $p$ of $D$ gains $e^{-\delta_1p/3}$, otherwise the string from an external vertex gives $e^{-\delta_2p}$; a vertex outside $F_v$ gives $e^{-\delta_2(4p/3-3)}$ by tree decay; a $\delta^{\rm im}$ line gives $e^{-\delta_0s_T/4}$; half of every exponent pays the vertex sums, whence $\delta'=\delta_2/6$, and $e^{-\delta'p(\delta_r)}\le C\delta_r^\sigma$ since $\mathrm p>1$; the Taylor remainders are $\le(Ce_rp_r^2)^{\bar n+1}\le C\delta_r^\sigma$. The universal-cover form replaces the periodic-cell reading of \cite[p.~337]{K2}: the periodised restriction of $\tilde A$ to a cell of side $\asymp p_r$ jumps by $\asymp p_r^2$ across the seam and is not regular, so the kernel bounds do not apply to it. (3) is as stated; the inversion symmetry about the coarse site used in the oddness rule requires the centred block convention for both flows (\Cref{rot:conv:C}(i)), and the frame-covariant average keeps it for the tiles.
\end{proof}

\begin{corollary}[the one-flow instance]\label{rot:cor:OF}
For one flow $F$ (standard or hybrid), a renormalised family of order $\le\bar n$ with lines from the slices $j=0,\dots,k$ of that flow (\cite[(3.76)]{K1}, each slice obeying \cite[(3.63)--(3.65)]{K1}; the coarse lines of depths $r$ and $r+n$ of the same flow included) and one element weakened by degree $\theta<\alpha_0$ obeys \cite[(1.33)]{Ba3} with an extra factor $C(1+k)^p$. Consequently King's estimate \cite[Prop.~3.6]{K1}, asserted on p.~665 without a printed argument, holds for both flows.
\end{corollary}
\begin{proof}
The multiscale bound on the point-split graph (the two summation points of a block product as vertices, the block kernel a line of index $r$): the slice lines satisfy the slice bounds with $d_l=-1$, a contour line of \cite[(3.26)--(3.27)]{K1} with $d_l=-1+\#(\text{contour ends})$ (\cite[(3.63)--(3.65), Prop.~3.7, Thm.~3.3]{K1} for the standard flow, \Cref{rot:prop:BD}(a) for the hybrid flow); the per-flow renormalised structure gives degree $\ge\alpha_0>0$ for every connected all-scale-summed subgraph except the trees \cite[(2.4)]{Ba3} and the counterterm-attachment subgraphs of \cite[(3.78)]{K1} with degree $0$, each of which costs a factor $(1+r)$ (King's ``possible addition of some power of $\ln(L^k\eps)^{-1}$'', p.~666, is $(1+k-r)^p$; both are $\le(1+k)^p$, \Cref{rot:lem:r}). The weakened line is the replacement $c_{l_0}\mapsto c_{l_0}\ell_j^{-\theta}$: the degree of every subgraph containing $l_0$ drops by $\theta<\alpha_0$ and stays positive, and no new zero-degree subgraph appears; a weakened vertex changes its exponent by $-\theta$; a weakened external field changes a norm only; the H\"older remainders of \cite[(3.83)]{K1} raise $d_l$ by $\alpha$ on the line they act on. This is King's own weakened-element induction \cite[(3.67)--(3.70)]{K1}, with no two-flow input; the strict positivity of the degrees after Ba{\l}aban's resummed rewriting is \cite[(3.5)]{Ba3}, and no cancellation for a fixed ordering of the lines is used.
\end{proof}

% r03_h1s_a: proof of the two-flow comparison bound (Theorem rot:thm:H1S), part 1: the all-C form, the slices, the elements, degrees per ordering, common moments, the current vertex.

\subsubsection*{Proof of the two-flow comparison bound}

The proof of \Cref{rot:thm:H1S} occupies the rest of this part. Its steps: the lines of every family are written as flow-independent lattice functions sliced by the heat kernel, plus coarse elements (this paragraph); degrees are read per ordering of the scale-summed elements, so that only the connected components of initial segments matter (\Cref{rot:h1s:MS}); each flow's expression is put in a normal form and one subtraction rule is applied, in each flow's own frame, to every component of degree $\le0$ that forms (\Cref{rot:h1s:rule}); after the rule every component containing the weakened element has degree $\ge\min(\alpha,1)$ (\Cref{rot:h1s:Dp}); every element type obeys a two-flow difference law with a positive exponent (\Cref{rot:h1s:P,rot:h1s:E,rot:h1s:U}); and the telescoping over the elements, together with the scale sums, gives the bound with the exponent $\theta_1$ (\Cref{rot:h1s:T} and the proof at the end).

\paragraph{The all-$C$ form and the slices.}
After the constant-background reduction (\Cref{rot:lem:CB}; at $e=0$ nothing is to be reduced) every line of a family is an $A=0$ line of its flow on the infinite lattice: $G^F_r=C^F-\Gamma^F$ (\Cref{rot:lem:Gam}(i)), where $C^F$ is the free propagator of the flow's lattice, $C^X(x,y)=C(x-y)$ and $C^Y(x'',y'')=C(R^{-1}(x''-y''))$ with \emph{one} lattice function $C=(-\Delta_{\eta'}+m_r^2)^{-1}$ on $\eta'\Z^3$ (for the photon, $D=(-\Delta_{\eta'}+\lmu^2\delta_r^2)^{-1}\otimes\mathbf 1_3$). Expanding every line of every graph in $C$ and $\Gamma$ gives \emph{$C/\Gamma$ monomials}, an exact identity; in each monomial the $C$-lines are the same lattice function evaluated in each flow's own frame, the vertices are matched by \Cref{rot:lem:vert} and the counterterm numbers are the same (\Cref{rot:lem:ctsame}), while $\Gamma^F$, $H^F$, $B^F=(A^F)^{-1}$ and the external fields depend on the flow. Write $C$ as the resolvent of the continuous-time simple random walk on $\eta'\Z^3$ with rate $\eta'^{-2}$ per direction, $C(u)=\eta'^{-3}\int_0^\infty e^{-m_r^2s}P_s(u)\,ds$, and slice it in the time variable:
\[
C^{(0)}\defeq\eta'^{-3}\!\int_0^{\ell_1^2},\qquad C^{(j)}\defeq\eta'^{-3}\!\int_{\ell_j^2}^{\ell_{j+1}^2}\ (1\le j\le r-1),\qquad C^{\rm top}\defeq\eta'^{-3}\!\int_1^\infty,
\]
all of $e^{-m_r^2s}P_s(u)\,ds$ ($\ell_r=1$). Gaussian and Poisson tail bounds for the walk give, for every $n\ge0$,
\begin{equation}\label{eq:h1s:B}
|\partial^nC^{(j)}(u)|\le c_n\ell_j^{-1-n}e^{-\delta|u|/\ell_j}\quad(0\le j\le r-1),\qquad|\partial^nC^{\rm top}(u)|\le c_ne^{-c'm_r|u|}(1+|u|)^{-1},
\end{equation}
$\partial$ a lattice difference; a derivative is charged to its line. The decomposition is the same for both flows as lattice functions; the $Y$ flow evaluates it at $R^{-1}(\tau x-\tau y)$, and in $X$-coordinates the $Y$-version obeys \eqref{eq:h1s:B} with $\delta/2$, since $|R^{-1}(\tau x-\tau y)|\ge|x-y|-2c_\tau\eta'$. The heat-kernel slices replace the block-spin slices of \Cref{rot:lem:slice} in the comparison of the all-$C$ lines; both decompositions are exact, and the heat-kernel one is flow-independent as a lattice function, which is what the subtraction rule below needs. (The gain $\eta'^{1/2}$ of \Cref{rot:cor:unif}(a) is recovered, and improved, by \Cref{rot:h1s:P}.)

\paragraph{Elements.}
The \emph{elements} of a monomial, with their per-flow bounds in the unit scale of level $r$, are: (a) the fine slice lines $C^{(j)}$, $D^{(j)}$, $0\le j\le r-1$, with their lattice derivatives, bounded by \eqref{eq:h1s:B}; (a$'$) the top pieces; (b) the coarse elements $\Gamma^F$, $H^F$, $B^F$ (\Cref{rot:lem:Gam}, \Cref{rot:lem:H}), bounded and decaying at the coarse mass, with one derivative and an $\alpha$-H\"older gradient (\cite[Thm.~3.3]{K1} for the standard flow, \Cref{rot:prop:BD}(a) for the hybrid flow); (c) the external fields $E^F\in\{H^F\phi_r,H^F_AA_r$, their covariant derivatives, the Cameron--Martin background$\}$, bounded by the field norms \cite[(3.2)]{K1} on $\supp\chi_r$ and $C^{1,\alpha}$ in the same sense; (d) the vertex functions $w$ (couplings, partitions of unity, the background), smooth with $\|\nabla w\|\le C$ (\Cref{rot:lem:vert}(ii)); (e) at $e>0$ the contour elements of \Cref{rot:prop:BD}(b): contour lines \cite[(3.27)]{K1}, double contour lines, contour vertices \cite[(3.24)--(3.26)]{K1} and the constant-background phases \cite[(3.46)]{K1}; (f) the counterterm numbers $\dm$, $E_1$, the same for both flows. Their matched differences, $X$-version at $x$ against $Y$-version at $\tau x$: (b) $|\Gamma^X-\Gamma^Y\circ\tau|\le C_\Gamma\eta^{1/2}$ (\Cref{rot:lem:Gam}(iii)), $|B^X-B^Y|\le0.2\eta^2$ (\Cref{rot:prop:fc}), coarse slices $\le C\eta^{1/2}\ell_j^{-1}e^{-\delta|x-y|/\ell_j}$ (\Cref{rot:cor:unif}(b)); (c) $|E^X(x)-E^Y(\tau x)|\le C\eta^{1/2}\|E\|$ (\Cref{rot:lem:H}, decaying form, applied to the lift); (d) $|w(x)-w(\tau x)|\le c_\tau\eta'\|\nabla w\|$; (e) contour lines: relative gain $(\eta'/\ell_j)^{1/2}$ on slice $j$, also with one lattice derivative; contour vertices: relative gain $\eta^{1/2}$ against their size $Ce_rp_r^2$; phases: $\le e\,p_r\lmu^{-1}c_\tau\delta_r^{-1/2}\eta'$ (\Cref{rot:prop:BD}(b)); the prefactor $p_r^2\delta_r^{-1/2}$ of the theorem comes from here and from nowhere else. Gradients of (b) and (c) are compared in \Cref{rot:h1s:E}.

\paragraph{The structures.}
With the vertices \cite[(3.18)--(3.23)]{K1} at $B=0$ (\Cref{rot:lem:CB}(1)), the connected all-$C$ line sets of non-positive degree that occur are the following (the list is proved in \Cref{rot:h1s:CL}). SE: the order-$e^2$ scalar self-energy, lines $\{C(x,y),D(x,y)\}$, two current vertices, both derivatives internal. VP: the order-$e^2$ vacuum polarisation, lines $\{C(x,y),C(x,y)\}$, two derivatives on the two lines, together with its seagull tadpole and derivative tadpole (\cite[(3.22)]{K1} with $m=2$), which enter only through the Ward sum of \Cref{rot:h1s:VP}. TAD: the $\lambda$-tadpole $C^{(j)}(0)$ and the seagull scalar tadpole $D^{(j)}(0)$, numbers, identical for both flows and cancelled exactly, slice by slice, by their counterterm shares (\cite[(3.33)]{K1}, \Cref{rot:lem:ctsame}). T: the order-$e^3$ triangle $\{C(a,b),C(b,c),D(a,c)\}$, three current vertices. SC: the order-$e^3$ seagull--current loop $\{C(s,b),D(s,b)\}$. W: the one-particle-irreducible two-point structures of order four: scalar (the $\lambda^2$ sunset; the $e^4$ rainbow, crossed, insertion and seagull graphs) and photon (the $e^4$ self-energy insertion, vertex correction, seagull graphs and the double seagull). L24: Ba{\l}aban's graph \cite[(2.4)]{Ba3}, a single scalar line between two bare current vertices carrying both their derivatives, $\partial\partial C$, with two external photon legs. CT: the counterterm share $c_\sigma$ of a self-energy or tadpole at slice assignment $\sigma$ (a one-site insertion, $|c_\sigma|\le C\ell_{\sigma_{\max}}^{-1}$, a loop line of dimension $-1$ at scale $\sigma_{\max}$), and the order-four share of dimension $0$. The $e^2\lambda$ two-point members are numbers or vanish identically; structures with an odd number of photon legs vanish by charge conjugation; the $(2,1)$-structure at order $e\lambda$ vanishes; four-leg loops have positive degree; contour vertices carry the explicit weight $\ge\frac12$ of \cite[(3.30)]{K1} at their site.

\paragraph{Degrees.}
For a connected set $H$ of elements (lines with their endpoints; a compound, defined in \Cref{rot:h1s:rule}, counts with the vertices and line dimensions it contains) put
\[
D(H)\defeq\sum_{l\in H}d_l+\sum_{v\in V(H)}a_v+3(|V(H)|-1),\qquad d_l=-1-(\text{number of lattice derivatives on }l),
\]
$a_v$ the vertex weight: an explicit factor $\eta'^a$, $a\ge0$, at $v$ counts as $\ell_{j(v)}^a$ with $j(v)\defeq\min\{j_l:l\ni v\}$; a derivative acting on an external field costs nothing, since the externals have bounded gradients. This is Ba{\l}aban's degree \cite[(2.1)--(2.3)]{Ba3} with every derivative charged to its line. The \emph{bare degree} $D_0(H)$ is the same count with the derivative placement of the normal form (NF2) below and before any derivative is transferred. All-$C$ degrees are integers; the only non-integer degrees are vertex weights $\alpha\in(0,1)$ produced by second-order subtractions with an external leg (\Cref{rot:h1s:U}(i)), so every degree is $n+m\alpha$ with $m\ge0$.

\begin{lemma}[the multiscale bound per ordering]\label{rot:h1s:MS}
Let $G$ be a connected set of elements of a monomial, each scale-summed element $l$ carrying the bound $c_l\ell_{j_l}^{d_l}e^{-\delta|x_l-y_l|/\ell_{j_l}}$ at its scale $j_l$ (fixed elements carry a bounded decaying kernel and impose no condition), and fix one ordering class of \cite[(2.7)]{Ba3}: an ordering $l(1),\dots,l(m)$ of the scale-summed elements with the scale sums restricted to $j_{l(1)}\le\dots\le j_{l(m)}$, ties assigned to one class. Suppose that every connected component of every initial segment $G_i\defeq\{l(1),\dots,l(i)\}$ has $D\ge0$, and $D\ge\theta$ if it contains a distinguished element $e_0$ whose constant carries the factor $(\eta'/\ell_{j(e_0)})^\theta$; let $n_0$ be the number of such components with $D=0$. Then the sum over the scales in the class and over all vertices but one is $\le C(1+r)^{n_0}\prod_lc_l$, and with the distinguished element $\le C\eta'^\theta(1+r)^{n_0}\prod_lc_l$, $C=C(G)$. Only the connected components of initial segments are tested.
\end{lemma}
\begin{proof}
This is the proof of \cite[Prop.~2.1]{Ba3} read for one ordering class, which is how Ba{\l}aban runs it: the finest element $l(1)$ is summed first, shrinking its line to a point with the exponent $D(G_1)$; the sum over $j_1\le j_2$ is geometric when $D(G_1)>0$ and costs a factor $(1+r)$ when $D(G_1)=0$; the merged vertex receives the weight $D(G_1)$, and the induction continues with $G_2,G_3,\dots$, at each stage testing only ``the zero-degree subgraphs met along the way'', the connected components of the initial segments. The same per-ordering positivity is King's \cite[p.~665]{K1}. With the distinguished element, every component containing it has exponent $D-\theta>0$ for the geometric sums, and the factor $\eta'^\theta\ell_{j(e_0)}^{-\theta}$ is absorbed by the merged vertex weight $\ell_{j'}^{D-\theta}\le\ell_{j(v)}^{D-\theta}$ at every later stage, so $\eta'^\theta$ survives to the end. A compound enters as a line of its stated dimension at the scale of its leg; its internal scale sums are done in \Cref{rot:h1s:U}(i) uniformly in the positions of its outer vertices, so the first step applies to it verbatim.
\end{proof}

\begin{lemma}[common moments]\label{rot:h1s:Z0}
Let $H(x,t)$ be a kernel built in a flow's own frame from $C$, $D$, their lattice derivatives, lattice Taylor remainders, lattice sums, the counterterm numbers and the bond shifts of \Cref{rot:h1s:cv} (no externals, no coarse elements). Then for $a=0,1$,
\[
\sum_{t\in\Lambda_X}H^Y(\tau x,\tau t)\,(\tau t-\tau x)^{\otimes a}=R^{\otimes a}\sum_{t\in\Lambda_X}H^X(x,t)\,(t-x)^{\otimes a},
\]
the indices compared frame-free (\Cref{rot:lem:vert}(iii)): zero moments are the same numbers and first moments the same vectors in both flows. An $\Oh$-invariant moment is isotropic: a rank-one or rank-three tensor vanishes, a rank-two one is $c\,\delta_{\mu\nu}$. On the tori the kernels are periodisations of the infinite-lattice ones and the moments agree up to $O(e^{-cm_r\cdot50MJ})$ (\Cref{rot:thm:V}(b)).
\end{lemma}
\begin{proof}
$H^Y$ is the conjugate of $H^X$ by the isometry $R$ (one lattice function evaluated in two frames) and $\tau:\Lambda_X\to\Lambda_Y$ is a bijection (\Cref{rot:def:match}), so the left side is $\sum_{v\in R\eta'\Z^3}H(R^{-1}v)v^{\otimes a}=R^{\otimes a}\sum_{u\in\eta'\Z^3}H(u)u^{\otimes a}$. The lattice model is $\Oh$-invariant (King's action on $\eps\Z^3$ with bond fields, in midpoint coordinates for bond kernels), so an $\Oh$-invariant tensor of odd rank vanishes under $-\mathbf 1\in\Oh$ and one of rank two is a multiple of $\delta_{\mu\nu}$.
\end{proof}

\begin{lemma}[the current vertex and the placement of its derivative]\label{rot:h1s:cv}
The current vertex \cite[(3.22)]{K1} ($m=1$) on the bond $b=(x,x+\eta'e_\mu)$ is $e_rA_b\,[\phi(b_+)\cdot q\phi(b_-)]/\eta'$ with $q$ antisymmetric. For two fields $a$, $b$ (a leg and an internal line after Wick contraction),
\[
(D_\mu a)(b)\cdot q\,b(b_-)=-a(b_+)\cdot q\,(D_\mu b)(b)+\eta'^{-1}\big[(a\cdot qb)(b_+)-(a\cdot qb)(b_-)\big],
\]
and summed against $A_b$ over the bonds the last term is $-\sum_x(\partial^*_\mu A_\mu)(x)\,(a\cdot qb)(x)$, the first term of \cite[(2.8)]{Ba3}; if the vertex function $g$ at the vertex is position dependent, the second term of \cite[(2.8)]{Ba3}, $\sum_x(a\cdot qb)(x)\sum_\mu(\partial^{\eta'}_\mu g)(x)A_\mu(x-\eta'e_\mu)$, appears as well. Hence the one lattice derivative of a current vertex may be written on either of its two scalar fields, the other being read at the shifted end of the bond, at the cost of a term in which the derivative acts on the vertex's photon field (and one with an explicit factor $\eta'$ from $\partial^{\eta'}g=\eta'\nabla g+O(\eta'^2)$). The identity is exact in each flow's own frame; a shifted evaluation point is matched by $\tau$ with $c_\tau$ replaced by $c_\tau+1$ in every matched-point bound. A current vertex has exactly two scalar fields and one photon field; the seagull $e_r^2|A|^2|\phi|^2$ and the quartic $\lambda_r|\phi|^4$ carry no derivative.
\end{lemma}
\begin{proof}
$a(b_+)=a(b_-)+\eta'(D_\mu a)(b)$ and $b(b_-)=b(b_+)-\eta'(D_\mu b)(b)$ give $a(b_+)\cdot qb(b_-)/\eta'=(D_\mu a)\cdot qb(b_-)+(a\cdot qb)(b_-)/\eta'=-a(b_+)\cdot q(D_\mu b)+(a\cdot qb)(b_+)/\eta'$; subtract, and sum by parts over the bonds $(x,x+\eta'e_\mu)$. For $a=b=\phi$ the remainder is $\phi\cdot q\phi/\eta'=0$, which is why the one-field form of the vertex has no such term; for two distinct fields it does not vanish, since the photon covariance is $\delta_{\mu\nu}$ times a scalar and not transverse.
\end{proof}

% r03_h1s_b: proof of the two-flow comparison bound, part 2: the sets of non-positive bare degree, the normal form, the subtraction rule.

\begin{lemma}[the connected all-$C$ sets of non-positive bare degree]\label{rot:h1s:CL}
Let $H$ be a connected set of all-$C$ lines with every current-vertex derivative placed on an internal line of its unit when the vertex belongs to a unit (the placement (NF2) below), and let $n_e$, $n_s$, $n_\lambda$ be the numbers of its current, seagull and quartic vertices, $n_{\rm ct}$ the number of its counterterm insertions, $E$ the number of its slots (legs of $V(H)$ not on lines of $H$, external fields included) and $x$ the number of current-vertex derivatives of $V(H)$ placed outside $H$. Then
\[
D_0(H)=\tfrac12n_e+n_s+n_\lambda+n_{\rm ct}+\tfrac12E-3+x .
\]
Decompose $H$ into its two-edge-connected blocks $P$ joined by $\rho-1$ bridges $b$. Then $D_0(H)=\sum_P[\omega(P)+\mathrm{out}(P)]+\sum_b(2-\mathrm{in}(b))\ge\rho-2$, where $\omega(P)$ is the degree of $P$ with every derivative of its vertices placed inside $P$, $\mathrm{out}(P)$ the number of derivatives of vertices of $P$ placed outside $P$ and $\mathrm{in}(b)\le2$ the number of derivatives on the bridge, so that $\sum_P\mathrm{out}(P)\ge\sum_b\mathrm{in}(b)$; and $\omega(P)\ge-1$ for every block with legs, with equality exactly for the units SE and VP (the tadpoles being dropped), for a counterterm insertion and, by convention, for a bare current vertex ($\omega=-1$, $\mathrm{out}=1$). Consequently $D_0(H)\le0$ forces one of: $\rho=1$ and $H$ is a unit ($-1$) or one of T, SC, W, the double seagull, a $\lambda^2$ or $e^4$ or $e^2\lambda$ two-point structure ($0$); $\rho=2$ and $H$ is two blocks with $\omega=-1$ and every non-internal derivative on the bridge: $U$--$l$--$U'$ ($l$ undifferentiated), $U$--$\partial l$--$v$, $v$--$\partial\partial l$--$v'$ ($=$ L24), and with a counterterm insertion ct--$l$--$U$, ct--$\partial l$--$v$, ct--$l$--ct; and, in the local-term monomials of the rule below, the same with a vertex of weight $-1$ (the vacuum-polarisation leftover of \Cref{rot:h1s:VP}) joined by an undifferentiated photon line to a vacuum-polarisation compound whose leg is that line, or to another such vertex. Unions of blocks at a cut vertex (a quartic or seagull vertex) have both slots at one vertex and are common numbers. Every member is a two-point structure except T and SC (three slots) and the one-point numbers.
\end{lemma}
\begin{proof}
Count: $\sum_ld_l=-I-n_\partial$ with $I$ the number of lines and $n_\partial$ the number of internal derivatives, $n_\partial=n_e-x$; the scalar legs give $2I_\phi+E_\phi=2n_e+2n_s+4n_\lambda+2n_{\rm ct}$ and the photon legs $2I_A+E_A=n_e+2n_s$; $|V|=n_e+n_s+n_\lambda+n_{\rm ct}$, and the counterterm insertions carry the vertex weights $\sum_va_v=-n_{\rm ct}$. Substituting in $-I-n_\partial+\sum_va_v+3(|V|-1)$ gives the formula. Blocks and bridges: $3(|V(H)|-1)=\sum_P3(|V(P)|-1)+3(\rho-1)$ because the block--bridge graph is a tree; a scalar bridge has $d_b=-1-\mathrm{in}(b)$ and a photon bridge carries no derivative in the placement (NF2); a derivative of a vertex of $P$ placed on a bridge counts $+1$ in $\mathrm{out}(P)$ and $-1$ on the bridge, and every derivative on a bridge is the own derivative of one of its end vertices, each of which lies in a block, whence $\sum_P\mathrm{out}(P)\ge\sum_b\mathrm{in}(b)$ and $D_0(H)\ge-\rho+2(\rho-1)=\rho-2$. The formula applied to a block with all derivatives internal gives $\omega(P)=\frac12n_e+n_s+n_\lambda+\frac12E-3$ for a block without insertions, and a counterterm insertion alone is a block with $E=2$ and $\omega=-1$; hence $\omega(P)\ge-1$ with the stated equality cases ($E=2$: $(n_e,n_s,n_\lambda)=(2,0,0)$; $E=3$: a single current vertex, whose derivative is necessarily outside), and the enumeration of the pairs $(P,P')$ with $\omega=-1$ and $\mathrm{in}(b)\le2$ gives the $\rho=2$ list; $\rho\ge3$ is impossible. The connected sets of bare degree $0$ with all derivatives internal are those with $\frac12n_e+n_s+n_\lambda+\frac12E=3$: $E=2$ with $(4,0,0)$, $(0,0,2)$, $(2,0,1)$, $(0,2,0)$, $(2,1,0)$, $(0,1,1)$; $E=3$ with $(3,0,0)$, $(1,1,0)$, $(1,0,1)$; $E=4$ with $(2,0,0)$, which is L24 (two current vertices joined by one scalar line carrying both derivatives; joined by a photon line both derivatives are outside and the degree is $2$). Integrality: $E\equiv n_e\pmod2$.
\end{proof}

\begin{corollary}[placement at the terminals; the derivative count]\label{rot:h1s:nd}
\textup{(i)} In a connected set with $D_0(H)\le0$ no vertex of $H$ places its own derivative outside $H$. \textup{(ii)} After the rule of \Cref{rot:h1s:rule} every line carries at most one lattice derivative at each end (its end vertex's own derivative or one transferred derivative, never both), hence at most two in all. The leg of a compound carries at most one as formed; when the compound is the differenced factor of (R2b) at one of its outer vertices, its leg carries the difference (\Cref{rot:h1s:U}(iv)) and then at most two; the middle line of a double compound carries none; a photon line carries only transferred derivatives, at most one per end; a contour line receives at most one, at its bond end.
\end{corollary}
\begin{proof}
(i) With $x\ge1$ the formula gives $D_0(H)\ge-1+x\ge0$, with equality only if the set with internal placement is a unit; but a unit's vertices have their derivatives inside by (NF2). (ii) Own derivatives: a current vertex has one, placed on one of its two scalar lines. Transferred derivatives land on a line $l$ at a vertex $w$ of the component $H$ being processed, $l\notin H$ later than $H$, $D(H)=0$ at formation, hence $D_0(H)\le0$ and, by (i), $w$'s own derivative is inside $H$: $l$ carries no own derivative of $w$. A second transfer onto $l$ at $w$ would come from a component $H'\ni w$ formed before $l$ enters the segment; $H'\supseteq H$ (nested initial segments), so $D(H')\ge D_0(H')+1\ge1$ unless $H'$ is a bare unit, which is never a component (\Cref{rot:h1s:Dp}(b)); so $H'$ is not processed. The leg of a compound is attached to the unit at a vertex whose own derivative is internal and at which no rule acts (the unit's lines are not elements before the compound enters, and afterwards they are inside it), so the leg carries at most one derivative, at its far end; the middle line of a double compound has two such ends. Photon lines carry no own derivative (\Cref{rot:h1s:cv}). For contour lines see (iii) of the rule.
\end{proof}

\begin{definition}[normal form and the subtraction rule]\label{rot:h1s:rule}
In each flow, in its own frame, before any element is mixed:
\begin{enumerate}[label=\textup{(NF\arabic*)},leftmargin=*]
\item Tadpoles and their counterterm shares cancel exactly per slice (TAD): they are dropped.
\item At every vertex of a unit (SE, VP, T, SC) whose derivative sits on a leg, the derivative is written on the internal scalar line by \Cref{rot:h1s:cv}; the remainder term, in which the derivative acts on the vertex's photon field, is the unit's own photon line, so it stays inside the unit (for SE the kernel is the sum of the pieces $\partial\partial C\cdot D$, $\partial C\cdot\partial D$, $C\cdot\partial\partial D$, each with $\sum_ld_l=-4$). After this a unit's leg carries at most the far vertex's derivative.
\item The seagull and derivative tadpoles of the order-$e^2$ vacuum polarisation are eliminated through the exact Ward sum of \Cref{rot:h1s:VP}; what remains at a vacuum-polarisation unit is the bubble for every slice assignment plus one common vertex weight.
\item The telescoping is done per ordering class of \cite[(2.7)]{Ba3}. The two classes that differ only by swapping the two internal lines of a vacuum-polarisation unit are merged before telescoping (identical element sets; \Cref{rot:h1s:MS} is applied per class in absolute value, a factor $2$ per such unit), so a vacuum-polarisation slice assignment is the unordered pair $\{j_1,j_2\}$ and its kernel the line-symmetrised one. The two Wick contractions of a self-energy between the same two legs lie in one class and are one element (their sum is the symmetrised kernel of \Cref{rot:h1s:M}).
\end{enumerate}
Fix an ordering of the scale-summed elements and process the initial segments $G_1\subset G_2\subset\dots$. Let $\sigma_{\max}(U)$ be the coarsest scale of a unit $U$.
\begin{enumerate}[label=\textup{(R\arabic*)},leftmargin=*]
\item\emph{Compounds.} A unit $U$ (SE, VP; order two) or a degree-zero one-loop piece (T, SC; order one) one of whose vertices $y$ has all its non-$U$ elements later than $\sigma_{\max}(U)$ (an external field counts as later) is bundled into the \emph{compound} $(U,F_y)$: the own-frame Taylor subtraction, about $U$'s other vertex $x$ (in leg-site coordinates, the vertex sites shifted by the bonds of the currents), of the product $F_y$ of the elements at $y$, taken on the earliest line $f$ of $F_y$ (or on the external if $F_y$ has no line), to second order for SE, VP and to first order for T, SC; one leg per unit, preferring a derivative-free, unclaimed leg. In the reduced ordering the compound is one element at the position of $f$ with the vertex set $V(U)\cup V(f)$ and the line dimensions of $U$ and $f$: $D(\{(U,f)\})=2+d_f\in\{1,0\}$ (SE, VP), $=3+d_f\ge1$ (T, SC); with an external leg $D=\alpha$ resp. $1$ (\Cref{rot:h1s:U}). Two units forced onto the same leg $l$ form the \emph{double compound} $(U_1,l,U_2)$, one element at $l$'s position with $D=D_0(U_1\cup l\cup U_2)=0$. The local terms of the subtraction are exact and taken in the flow's own frame: for SE the zeroth-order term $c_\sigma f(x)g(x)$ cancels the co-switched counterterm share of the same $\sigma$ (\cite[(3.33)]{K1} through order four; \Cref{rot:lem:ctsame}) and the first-order term vanishes by \Cref{rot:h1s:M}; for VP the zeroth-order term enters the Ward sum of \Cref{rot:h1s:VP} and the first-order term vanishes by \Cref{rot:h1s:M}; for T the local term is $f(x)\sum_zT''(x,z)g(z)$ with the two-point kernel $T''(x,z)\defeq\sum_yT(x;y,z)$ of dimension $-3$, a degree-zero line processed by (R2); for SC the local term vanishes or carries an explicit factor $\eta'$ (the bond shift). A unit none of whose vertices qualifies (both legs earlier than $\sigma_{\max}$, the \emph{split} unit) is not an element: its lines are ordinary elements and its counterterm share stands alone as a loop line of dimension $-1$ at scale $\sigma_{\max}$ in a separate monomial.
\item\emph{Components of degree zero.} Whenever a newly formed connected component $H$ of the current initial segment has $D(H)=0$ (degrees after the earlier transfers), apply the first available clause.
\begin{enumerate}[label=\textup{(\alph*)},leftmargin=*]
\item If $H$ contains a bare current vertex $w$, or a vertex $w$ of a split unit, whose derivative sits on a line of $H$, whose other scalar field is not in $H$ and whose photon field is an external field or a photon line not in $H$: rewrite the vertex by \Cref{rot:h1s:cv}. This gives two monomials, (i) with the derivative on the other scalar field and (ii) with $\partial^*$ on the photon field; in both $D(H)$ rises to $1$, the recipient (later than $H$, since $H$ is a component) changing $d\mapsto d-1$, or at no cost if it is an external field; the term with $\partial^{\eta'}g$ carries an explicit $\eta'$, a vertex weight $1$. Each vertex is rewritten at most once. If $w$'s photon field is a line inside $H$, the clause is not available (the move would leave $D(H)=0$): this happens exactly for the outer vertices of the two-point structures W and for the leg-carrying vertices of a triangle, which go to (b) and (R1).
\item Else, if $H$ is a two-point structure (exactly two vertices $x$, $t$ of $H$ carry elements outside $H$ or externals): the first-order subtraction at the terminal $t$ about $x$ (either terminal may be chosen),
\[
\textstyle\sum_tH(x,t)G(t)=c\,G(x)+\sum_tH(x,t)[G(t)-G(x)],\qquad c\defeq\sum_tH(x,t),
\]
$G$ the product of the later elements at $t$, Leibniz-expanded so that exactly one factor is differenced. $c$ is a common number of both flows (\Cref{rot:h1s:Z0}), bounded by \Cref{rot:h1s:MS} applied per flow to $H$ with one vertex fixed ($|\sum_\sigma c_\sigma|\le C(1+r)^{n_0+1}$); for the sunset and the scalar $e^4$ structures it is $0$ exactly with the co-switched order-four counterterm share, for the $e^4$ vacuum polarisation it is $c\,\delta_{\mu\nu}$, for a triangle kernel $T''$ it is $0$ exactly (the summand is odd under the lattice inversion through a site: the current on a bond goes to minus the current on the inverted bond, the propagators are even and the three lines are distinguishable), and for a double compound standing alone it is $0$ exactly (translation invariance of the outer unit's kernel). $D(H)$ rises to $1$ (the factor $|t-x|$ absorbed on a line of $H$); the differenced factor changes $d\mapsto d-1$ if it is a line (for a compound, its leg, \Cref{rot:h1s:U}(iv)), at no cost if it is an external, a coarse element or a vertex function; the $Y$-displacement $|\tau t-\tau x-(t-x)|\le2c_\tau\eta'$ against the $Y$-gradient is an explicit-$\eta'$ vertex weight, a degree-zero component with gain $\eta'\ge\eta'^\theta$, counted among the logarithms.
\item Else $H$ has at most one vertex with outside elements or externals and is a common number (\Cref{rot:h1s:Z0} with $a=0$; both flows evaluate the same closed lattice sum), bounded per flow by \Cref{rot:h1s:MS}: a vertex weight $0$, never a switched element.
\end{enumerate}
\item\emph{Local-term monomials.} The monomials produced by the zeroth-order terms of (R1) and (R2b) (the unit or structure collapsed to a vertex of weight $0$, or $-1$ for the vacuum-polarisation leftover of \Cref{rot:h1s:VP}) are processed by the same rules; their degree-zero components are L24-type lines, the weight classes of \Cref{rot:h1s:CL}, and, when a triangle inside a W is collapsed, the kernels $\{\partial C,D\}$ and $\{C,\partial C\}$ between the collapsed point and the other terminal, all handled by (R2b) with common numbers.
\end{enumerate}
\textup{(iii)} At $e>0$ the recipient of (R2a)(ii) may be a contour line \cite[(3.27)]{K1} whose bond end sits at $w$; it receives this one derivative at its bond end and no other, since its contour end is the site of a contour vertex, which carries the weight $\ge\frac12$ and so never lies in a degree-zero component. Contour vertices are never recipients of any clause, for the same reason, and the constant-background phases \cite[(3.46)]{K1} are not photon fields. The lifted external background at $w$ is an external field with the gradient bound $\le Cp_r$ (\Cref{rot:lem:CB}(0)(R)).
\end{definition}

The rule is Ba{\l}aban's own device for the graph (2.4), ``the first few $G_1,G_2^{(2)},\dots$ are graphs (2.4)\dots replaced by graphs with degree $+1$'' by the integration by parts \cite[(2.8)]{Ba3}, applied to every degree-zero component that forms in an initial segment, with composite lines (compounds) allowed and with the first-order subtraction (R2b) for the components without a usable bare vertex. Two features are forced by the two flows: the units are bundled with a leg as one element, so that no line of a unit is ever processed alone inside an isolated unit (processing the lone $\partial\partial C$ of a self-energy with both legs later would lose the gain, since the Taylor order at the slot would be fixed before the second unit line enters), and the local terms are common numbers of both flows, so that no first moment of a flow-dependent object is ever mixed between the frames.

% r03_h1s_c: proof of the two-flow comparison bound, part 3: degree bookkeeping per ordering after the rule; the enumeration.

\begin{lemma}[degree bookkeeping]\label{rot:h1s:Dp}
Fix an all-$C$ component $G$ of a monomial, a derivative placement, an ordering of its scale-summed elements, and apply (NF1)--(NF4) and (R1)--(R3). Then in the reduced ordering every connected component $H$ of every initial segment satisfies
\[
D(H)\ge1,\quad\text{or}\quad D(H)=n+m\alpha\ \ (n\ge0,\ m\ge1;\ H\text{ contains a compound on an external leg}),\quad\text{or}\quad H\text{ is an explicit-$\eta'$ component}
\]
(a $Y$-displacement term of (R2b) or a $\partial^{\eta'}g$ term of (R2a), with $D(H)=1-\theta$ after the gain $\eta'\ge\eta'^\theta$ is taken out), with at most $n_0\le\bar n$ zero-degree components in the per-flow count (local vertices of weight $0$, explicit-$\eta'$ components). In particular $D(H)-\theta\ge\min(\alpha,1)-\theta>0$ for $\theta\le\theta_1$ whenever $H$ contains the weakened element.
\end{lemma}
\begin{proof}
(a) Degrees are a property of the multilinear form and flow-independent: both flows have the same lines, the same placements and the same transfers, since (NF2) and the rule are applied in each flow's own frame to the same combinatorial objects.

(b) \emph{No component has $D<0$.} By \Cref{rot:h1s:CL} a connected line set with $D_0<0$ is a unit $U$. In the reduced ordering $U$'s lines are either inside a compound, and then $U$ alone is not an element, or $U$ is split, and then both legs $f$, $g$ of $U$ are earlier than $\sigma_{\max}(U)$: the first initial segment containing both lines of $U$ contains $f$ and $g$, attached at $U$'s vertices, so the component is $\{f,g,U,\dots\}$ with $D\ge d_f+d_g-4+9\ge1$ ($d_f,d_g\ge-2$ by \Cref{rot:h1s:nd}(i)); the segment containing only $U$'s first line, the lone $\partial\partial C$ of a split self-energy with $D=0$ when it precedes both legs, is a degree-zero component covered by (c). Transfers never create $D<0$: a transfer lands on an element later than its source component $H_0$ and attached to a vertex of $H_0$, so every later component containing the recipient contains $H_0$, where the $+1$ and the $-1$ cancel (this is \cite[(2.8)--(2.9)]{Ba3}: ``the degrees of these graphs are $\ge D(G_0)$''); the components containing $H_0$ but not the recipient rose by $1$; no other component changed. Loop closure after a transfer is harmless for the same reason: if the recipient line later closes a loop, the component $H\cup l$ has $D(H\cup l)\ge D_0(H\cup l)\ge-1$, with $-1$ only for a bare unit; then $H$ is the first line of a split unit and $l$ its other line, the coarsest, so both legs are earlier and attached when $l$ enters, and the component is $\{f,g,U,\dots\}$ with $D\ge1$ as above.

(c) \emph{Components with $D(H)=0$ at formation.} Since every transfer inside $H$ with recipient inside $H$ is neutral and one with recipient outside is $+1$, $D(H)\ge D_0(H)$, so $D_0(H)\in\{-1,0\}$; $-1$ is excluded by (b), so $H$'s line set is in the list of \Cref{rot:h1s:CL} (compounds expanded; for a derived structure the degree is that of its line set). Case by case:
\begin{itemize}[nosep,leftmargin=*]
\item $H=(U,\partial f)$ with the far vertex $w$ of $f$ bare, or $H=$ L24, or $H=\{\partial f,{\rm ct}_\sigma\}$ (a split form), or $H=(U,\partial f)$ with the derivative transferred onto $f$ earlier: a bare vertex $w\in H$ has its derivative on a line of $H$, its other scalar field outside $H$ and its photon field outside $H$ (external, or a photon line later than $H$); (R2a) raises $D$ to $1$. The lone $\partial\partial C$ of a split self-energy: (R2a) at a vertex of the split unit (its photon is the unit's $D$ line, outside the one-line component), or, equivalently, (R2b) with $c=\sum_y\partial\partial C(x,y)=0$ exactly (a lattice total derivative), the recipient being the leg at a terminal, which precedes the second unit line in the split regime; then $\{\partial\partial C^{+1},\nabla f\}$ has $D=2$ and $\{\partial\partial C^{+1},\nabla f,D\}$ has $D=1$.
\item $H=T''$ (the triangle's two-point local kernel): the terminals' photon field is the triangle's own $D$ line, inside $T''$, so (R2a) is not available; (R2b) applies with $c=0$ exactly (inversion oddness); $D\mapsto1$.
\item $H=$ W, scalar (outer current vertices bare, derivatives internal): an outer vertex's photon is a line inside W; (R2b) with $c=0$ exactly against the co-switched order-four counterterm share; $D\mapsto1$. $H=$ W, photon, or the double seagull, the sunset, the $e^2\lambda$ members: (R2b) with the common number $c$; $D\mapsto1$. When two triangles inside a W are both bundled (the crossed and vertex-correction graphs: each triangle's outer vertex carries an external leg or a later line, so each qualifies under (R1)), their line sets overlap and the two compounds merge into one element whose line set is all of W: a two-point degree-zero component when it stands alone, handled by (R2b). The collapsed local terms inside W, $\{\partial C,D\}$ and $\{C,\partial C\}$ between the collapsed point and the other terminal, are degree-zero two-point kernels of the SC and VP type: (R2b) with the common number; $D\mapsto1$.
\item Chains $U_1$--$l$--$U_2$, ct--$l$--$U$, ct--$l$--ct with $D=0$: each $U_i$ is bundled on the chain line $l$ or on its outer leg (if neither vertex of $U_i$ qualifies, its outer leg is earlier and attached, and then $H$ contains it with $D\ge1$); if $U_1$ bundles on $l$ and $U_2$ on its outer leg $g$, $\{(U_1,l),(U_2,g)\}$ has $D=(2+d_l)+(2+d_g)\ge1$; if both bundle on $l$, $H$ is the double compound element, which with two units is never alone (its outer legs are earlier) and in chains of three or more units can stand alone (the middle unit forced onto a leg already claimed by a neighbour while both outer legs of the double compound are later): then it is a two-point degree-zero component with no bare vertex and (R2b) applies, with $c=0$ exactly; $D\mapsto1$. A counterterm share stands alone only when its unit is split, so $\{f,g,{\rm ct}\}\ge5+d_f+d_g\ge1$, and $\{\partial f,{\rm ct}\}$ alone never forms since $g$ is earlier too.
\item $H$ one-point (both slots at one vertex: the seagull with two external photons and a hanging self-energy; the all-$C$ components of vacuum monomials cut at their $\Gamma$-lines): a common number, (R2c).
\item The weight classes of the local-term monomials, a weight-$(-1)$ vertex joined by an undifferentiated photon line to a vacuum-polarisation compound whose leg is that line, or to another such vertex ($-1+(-4-1)+3\cdot2=0$, $-1-1-1+3=0$): two-point structures whose outer vertices carry external or later photon fields; (R2b) with the common number; $D\mapsto1$.
\end{itemize}
Nothing else occurs: the list of \Cref{rot:h1s:CL} is exhausted (T and SC are bundled when a vertex qualifies, and otherwise never alone: $\{f_{\rm earlier},T\}$ has $D=d_f+3\ge1$).

(d) \emph{Components with $D=\alpha$.} A compound $(U,E)$ on an external leg alone; any component containing it and a line has $D\ge\alpha+1$.

(e) \emph{Termination and the count.} Each application of (R2) raises one freshly formed component and rewrites one vertex or differences one factor; a vertex is rewritten at most once (after (R2a) at $w$ every later component containing the recipient contains $H_0$ with $D\ge1$); (R2b) is applied once per formed two-point component; the local-term monomials have fewer vertices. The per-flow zero-degree components are the weight-$0$ local vertices and the explicit-$\eta'$ components, at most one per rule application, so $n_0\le\bar n$.
\end{proof}

\begin{remark}[the enumeration]\label{rot:h1s:census}
The case analysis of \Cref{rot:h1s:Dp}(c) is a finite statement for each graph, and an exhaustive enumeration, in exact integer arithmetic, confirms it for every graph with at most two units among the two- and three-point structures through order four (the isolated self-energy and vacuum polarisation, the tree (2.4), the triangle and the seagull--current loop with their leg currents, a unit between two currents, the rainbow, crossed, insertion, vertex-correction and seagull graphs, the chains $U$--$l$--$U$ and $U$--$\partial l$--$w$--$l'$--$U$ with bare and unit vertices, the collapsed local-term graphs and the weight classes): for every derivative placement allowed by \Cref{rot:h1s:cv} and every ordering of the scale-summed lines, walking the initial segments and applying the rule as stated (both monomials of (R2a) followed separately), no component of negative degree forms, every degree-zero component falls under one of the clauses, the rule terminates, and no line carries more than two lattice derivatives; the same holds for the chains of three and four units, the rings and the counterterm monomials, the double compound standing alone in a fraction of the orderings of the three-unit chain and being handled by (R2b). The proof for every family of order $\le\bar n$ is \Cref{rot:h1s:CL} together with (b)--(e), which use only the block arithmetic and the nesting of the initial segments. The statement for chains of any number of units is degree arithmetic: a third unit cannot join a set of bare degree $\le0$ (\Cref{rot:h1s:CL}, $\rho\le2$), and each further unit enters as a compound at a later leg or inside a component of degree $\ge1$.
\end{remark}

% r03_h1s_d: proof of the two-flow comparison bound, part 4: first moments of the units, the vacuum polarisation, the slice law, gradients of externals, the compound as a weakened element.

\begin{lemma}[the first moments of the units vanish]\label{rot:h1s:M}
In each flow's own frame: \textup{(i)} for every slice assignment $\sigma$ the self-energy kernel $k_\sigma$, as it enters any family (the sum of the two Wick contractions between the same two legs, in leg-site coordinates), has $F_\sigma\defeq\sum_uk_\sigma(u)\,u=0$; \textup{(ii)} for every unordered slice assignment $\sigma=\{j_1,j_2\}$ of the vacuum polarisation, the line-symmetrised kernel $\Pi_{\sigma,\mu\nu}\defeq\Pi_{(j_1,j_2),\mu\nu}+\Pi_{(j_2,j_1),\mu\nu}$ has $F_{\sigma,\mu\nu\rho}\defeq\sum_y\Pi_{\sigma,\mu\nu}(x,y)(y-x)_\rho=0$, in site and in bond-midpoint coordinates alike. (For a labelled assignment $(j_1,j_2)$ with $j_1\ne j_2$ the first moment is not zero; this is why (NF4) merges the two classes.)
\end{lemma}
\begin{proof}
Why it is needed: the compound's expansion $\sum_yk_\sigma(x,y)f(y)=c_\sigma f(x)+F_\sigma\cdot\nabla f(x)+\sum_yk_\sigma(x,y)T_xf(y)$ removes the first-order term. For a slice-line leg a term $F_\sigma\cdot\nabla f(x)$ would be harmless even if $F_\sigma\ne0$ ($F^Y=RF^X$ by \Cref{rot:h1s:Z0}). For an external or coarse leg it would not: $F^X\cdot\nabla E^X(x)$ and $F^X\cdot R^{-1}\nabla E^Y(\tau x)$ differ by $F^X\cdot(1-R^{-1})\nabla E(x)$ up to small terms, a frame term of order $|F|\|E\|$ with no gain in $\eta'$. (i) After (NF2) the self-energy at slice $\sigma$ is a bilinear form $B_\sigma(\phi,\phi)=\sum_{x',y'}\phi(x')\cdot\tilde k_\sigma(x',y')\phi(y')$ in the two leg fields, $x'$, $y'$ the leg sites (vertex sites shifted by the bonds of the currents), with $\tilde k_\sigma$ proportional to the identity in the $O(2)$ indices ($q^2=-1$; no $q$-part by charge conjugation). In every family the self-energy sits between two leg elements $g$ at $x'$ and $f$ at $y'$ through both Wick contractions, which have identical element sets; their sum is $\sum g(x')[\tilde k_\sigma(x',y')+\tilde k_\sigma(y',x')]f(y')$. The kernel $k_\sigma\defeq\tilde k_\sigma+\tilde k_\sigma^T$ is translation invariant and symmetric, so $\kappa_\sigma(u)\defeq k_\sigma(x',x'+u)$ is even and $F_\sigma=0$, whatever the rewriting of \Cref{rot:h1s:cv} did inside the kernel; the expansion is about the leg site, so no shifted-site moment appears. (ii) The bubble at labelled slices is $\Pi_{(j_1,j_2),\mu\nu}(x,y)\propto C^{(j_1)}(x+e_\mu,y+e_\nu)C^{(j_2)}(x,y)-C^{(j_1)}(x+e_\mu,y)C^{(j_2)}(x,y+e_\nu)$; under $-\mathbf 1\in\Oh$ the two factor positions swap slices, so the labelled kernel is $\Oh$-covariant only up to $j_1\leftrightarrow j_2$ and the symmetrised one is the $\Oh$-covariant object, a kernel on bonds covariant in midpoint coordinates, $\Pi_\sigma(Rm,Rm')=R\Pi_\sigma(m,m')R^T$ for $R\in\Oh$. Its midpoint first moment is an $\Oh$-invariant rank-three tensor, hence $0$; its zero moment is $c_\sigma\delta_{\mu\nu}$. In site coordinates $m'-m=(y-x)+(e_\nu-e_\mu)\eta'/2$, so $F_{\sigma,\mu\nu\rho}=0-c_\sigma\delta_{\mu\nu}(e_\nu-e_\mu)_\rho\eta'/2=0$.
\end{proof}

\begin{lemma}[the vacuum polarisation]\label{rot:h1s:VP}
\textup{(a)} With every line sliced, the order-$e^2$ photon two-point function at zero momentum vanishes in each flow, \cite[(2.23)--(2.28)]{Ba3} (valid with free boundary conditions) and \cite[\S3.5]{K1}: $\sum_{\sigma}c_{\sigma,\mu\nu}+\sum_jc^{\rm sg}_j\delta_{\mu\nu}+c^{\rm dt}_{\mu\nu}+c^{\rm top}_{\mu\nu}=0$, where $c_\sigma\defeq\sum_y\Pi_\sigma(x,y)$ is the bubble's zero moment, $c^{\rm sg}_j$ the seagull tadpole, $c^{\rm dt}$ the derivative tadpole and $c^{\rm top}$ the top pieces, all common numbers, all sums finite ($\sum_\sigma|c_\sigma|\le C\eta'^{-1}$ at fixed $r$). \textup{(b)} The seagull and derivative tadpoles are single-vertex structures of degree $-1$ and may not appear as elements; (a) replaces them by $-\sum_\sigma c_\sigma-c^{\rm top}$. With the local terms $+c_\sigma f(x)g(x)$ of the bundled bubbles ((R1); $f$, $g$ the photon legs) the complete local content at a vacuum-polarisation vertex is $[-\sum_{\sigma\ \rm unbundled}c_\sigma-c^{\rm top}]f(x)g(x)$, a slice assignment being unbundled iff $\sigma_{\max}>J\defeq\max(j_f,j_g)$, plus the remainders of the bundled bubbles and the unbundled bubbles themselves, which are never alone ($\{f,\Pi_\sigma\}$ has $D=1$). Since $|c_\sigma|\le C\ell_{\sigma_{\min}}/\ell_{\sigma_{\max}}^2$ (placement $(1,1)$) and $\le C\ell_{\sigma_{\min}}^2/\ell_{\sigma_{\max}}^3$ (placement $(2,0)$, using $\sum_u\partial\partial C^{(j_1)}(u)=0$),
\[
\Big|\sum_{\sigma_{\max}>J}c_\sigma\Big|\le C\sum_{j_2>J}\ell_{j_2}^{-2}\sum_{j_1\le j_2}\ell_{j_1}\le C\ell_J^{-1},\qquad|c^{\rm top}|\le C:
\]
a vertex weight $-1$ at scale $J\ge j(x)$ plus a weight $0$. Photon legs carry no own derivative, so $\{f,\text{weight}\}$ has $D=1$ and $\{f,g,\text{weight}\}$ has $D=3$. The Ward identity is used once, per flow, for the sum over all $\sigma$; its fine end, $\sum_{\sigma_{\max}\le J}c_\sigma$, of order $\eta'^{-1}$, is exactly what the bundled remainders carry, and only the coarse end survives as a weight. With external photon legs the leftover is $0$ exactly. \textup{(c)} The order-$e^4$ vacuum polarisation, and every degree-zero two-point structure handled by (R2b), needs no identity: $c_{\mu\nu}\defeq\sum_t\Pi^{(4)}_W(x,t)$ is a common number (\Cref{rot:h1s:Z0}), isotropic, and $|\sum_\sigma c_\sigma|\le C(1+r)^{n_0+1}$ by \Cref{rot:h1s:MS} applied per flow to W with one vertex fixed (all proper components of W are already of degree $\ge0$ after the rule; $D(W)=0$ gives one logarithm). The local term $c\,A^F(x)\cdot A^F(x)$ is a weight-$0$ vertex term common to both flows; its difference is telescoped over the externals. \textup{(d)} In the scalar sector the bundled self-energy's zeroth-order term cancels against ${\rm ct}_\sigma$ exactly; the sunset and the scalar $e^4$ structures give $c=0$ exactly with the co-switched order-four share.
\end{lemma}

\begin{lemma}[the slice law]\label{rot:h1s:P}
Let $g_j\defeq C^{(j)}-C^{(j)}_c$ with $C^{(j)}_c(v)\defeq\int_{\ell_j^2}^{\ell_{j+1}^2}e^{-m_r^2s}(4\pi s)^{-3/2}e^{-|v|^2/4s}ds$ the isotropic continuum slice, $3\le j\le r-1$. Then for every $n\ge0$ lattice derivatives
\[
|\partial^ng_j(u)|\le C_n(\eta'/\ell_j)^{1/2^n}\ell_j^{-1-n}e^{-\delta_n|u|/\ell_j},
\]
and consequently, for matched points $|u'-R^{-1}u|\le4c_\tau\eta'$ (bond shifts included),
\[
|\partial^nC^{(j)}(u)-R^{\otimes n}\partial^nC^{(j)}(u')|\le C(\eta'/\ell_j)^{1/2^n}\ell_j^{-1-n}e^{-\delta_n|u|/(2\ell_j)},
\]
the same for the photon slices and, with $\ell_j$ replaced by $1$ and the mass decay, for the top piece. For $j\le2$ there is no gain and the factor $(\eta'/\ell_j)^\theta$ is paid by a degree loss $\theta$ (\Cref{rot:cor:unif}(a), King's device at the finest slices \cite[p.~675]{K1}). Constants are uniform in $m_r\in(0,1]$ and $r$.
\end{lemma}
\begin{proof}
Value clause $n=0$: the coordinates of the continuous-time walk are independent one-dimensional walks \cite[Prop.~1.2.2]{LL10}, and the local limit law \cite[Thm.~2.5.6]{LL10} gives, for $|u|\le s/(8C_0\eta')$ and $s\ge\ell_3^2$, $|\eta'^{-3}P_s(u)-(4\pi s)^{-3/2}e^{-|u|^2/4s}|\le c(\eta'/\sqrt s)(4\pi s)^{-3/2}e^{-|u|^2/8s}$; integrated over $s\in[\ell_j^2,\ell_{j+1}^2]$ this is $\le c(\eta'/\ell_j)\ell_j^{-1}e^{-\delta|u|/\ell_j}$, and beyond that range both kernels are bounded by their Gaussian and Poisson tails, $\le C_N(\eta'/\ell_j)^Ne^{-\delta|u|/\ell_j}$; the geometric mean with the decay of each kernel halves the rate and keeps the factor. Derivatives by interpolation: for a lattice function $h$ and $\rho=N\eta'$, $\partial_\mu h(u)=[h(u+\rho e_\mu)-h(u)]/\rho-\rho^{-1}\sum_{i=1}^{N-1}(\rho-i\eta')\partial_\mu\partial_\mu h(u+i\eta'e_\mu)\eta'$, so $|\partial_\mu h(u)|\le2\sup|h|/\rho+\rho\sup|\partial\partial h|$ over the segment. With $h=\partial^ng_j$, the crude bound $|\partial^{n+2}g_j|\le|\partial^{n+2}C^{(j)}|+|\partial^{n+2}C_c^{(j)}|\le C\ell_j^{-3-n}e^{-\delta|u|/\ell_j}$ from \eqref{eq:h1s:B} and the smoothness of $C_c$, and $\rho\defeq\ell_j\epsilon_n^{1/2}$ with $\epsilon_n\defeq C_n(\eta'/\ell_j)^{1/2^n}$ (so $\rho\ge\eta'$; round to a multiple of $\eta'$), one gets $|\partial^{n+1}g_j|\le C\epsilon_n^{1/2}\ell_j^{-2-n}$, the decay rate kept since the segment has length $\le\ell_j$; induction on $n$. Matched points: $\partial^nC^{(j)}(u)-R^{\otimes n}\partial^nC^{(j)}(u')=\partial^ng_j(u)+[\partial^nC_c^{(j)}(u)-R^{\otimes n}\partial^nC^{(j)}_c(u')]-R^{\otimes n}\partial^ng_j(u')$; the middle term is the rotation covariance of the isotropic $C_c$, $\nabla^nC_c(Rv)=R^{\otimes n}\nabla^nC_c(v)$, plus the displacement $|R^{-1}u-u'|\le4c_\tau\eta'$ against $|\nabla^{n+1}C_c|\le C\ell_j^{-2-n}$, and the lattice-difference-versus-derivative error of $C_c$, $\le C\eta'\ell_j^{-2-n}$; both are inside the claim. (The true law is $(\eta'/\ell_j)^2$ for the values and $\eta'/\ell_j$ for first differences; the exponents $1/2^n$ are far from sharp and suffice, since every gain is compared with $\theta_1\le\frac1{16}$.)
\end{proof}

\begin{lemma}[gradients of externals and coarse elements]\label{rot:h1s:E}
For $E\in\{H^F\phi_r,H^F_AA_r$, their covariant derivatives, the background$\}$ and for $\Gamma^F(\cdot,z)$, $H^F(\cdot,z)$: $|E^X(x)-E^Y(\tau x)|\le C\eta^{1/2}\|E\|$ (\Cref{rot:lem:H}, \Cref{rot:lem:Gam}(iii)), and both are restrictions of $C^{1,\alpha}$ functions $g^X$, $g^Y$ on $\R^3$ with $\|\nabla g^F\|_\infty\le C_\nabla\|E\|$ and $[\nabla g^F]_\alpha\le C\|E\|$ (\cite[Thm.~3.3, (3.8)]{K1} for the standard flow, \Cref{rot:prop:BD}(a) for the hybrid flow; for the lifts, the Fourier representation of \Cref{rot:lem:H} with smooth symbols). Then, with $\kappa\defeq\alpha/(2(1+\alpha))$, $|\nabla g^X(x)-\nabla g^Y(x)|\le C\eta^\kappa\|E\|$ for all $x\in\R^3$, in particular at $x$ and at $\tau x$. Consequently the second-order own-frame Taylor remainders $T^F_xE^F(y)\defeq E^F(y)-E^F(x)-(y-x)\cdot\nabla g^F(x)$ (in the $Y$ frame at $\tau x$, $\tau y$) differ by
\[
|T^X_xE^X(y)-T^Y_{\tau x}E^Y(\tau y)|\le C\|E\|\min\{|u|^{1+\alpha}+\eta'^{1+\alpha},\ \eta^{1/2}+|u|\eta^\kappa+\eta'\},\qquad u=y-x .
\]
For a fine slice line $f=\partial^aC^{(j_f)}(\cdot-z)$ in place of $E$ the same holds with $\eta^{1/2}$, $\eta^\kappa$, $\eta'$ replaced by $(\eta'/\ell_{j_f})^{1/2^a}\ell_{j_f}^{-1-a}$, $(\eta'/\ell_{j_f})^{1/2^{a+1}}\ell_{j_f}^{-2-a}$, $\eta'\ell_{j_f}^{-2-a}$, times $e^{-\delta\,{\rm dist}/(2\ell_{j_f})}$, and with the remainder bound $c|u|^{1+\alpha}\ell_{j_f}^{-2-a-\alpha}$ (\Cref{rot:h1s:P}; no interpolation needed).
\end{lemma}
\begin{proof}
$h\defeq g^X-g^Y$ has $\|h\|_\infty\le C\eta^{1/2}\|E\|$ (at lattice points by the value difference and $|g^Y(\tau x)-g^Y(x)|\le C_\nabla c_\tau\eta'\|E\|$, between them by the Lipschitz bound) and $[\nabla h]_\alpha\le2C\|E\|$; for a unit vector $e$ and $\rho>0$, $|e\cdot\nabla h(x)|\le|h(x+\rho e)-h(x)|/\rho+[\nabla h]_\alpha\rho^\alpha\le2C\eta^{1/2}\|E\|/\rho+2C\|E\|\rho^\alpha$; choose $\rho=\eta^{1/(2(1+\alpha))}$. The first remainder bound is the sum of the two remainders ($|\tau y-\tau x|\le|u|+2c_\tau\eta'$); the second is the value law at $y$ and $x$, the gradient law at $x$ and the displacement $|(\tau y-\tau x)-(y-x)|\le2c_\tau\eta'$ against $\|\nabla g^Y\|$.
\end{proof}

\begin{lemma}[the compound as a weakened element]\label{rot:h1s:U}
Let $B=(U,f)$ be a compound of (R1): $U$ a self-energy or vacuum polarisation with slice assignment $\sigma$ between the leg sites $x$ and $y$, placement-summed and symmetrised kernel $k_\sigma(x,y)=\kappa_\sigma(x-y)$ with $|\kappa_\sigma(u)|\le C\ell_{j_1}^{d_1}\ell_{j_2}^{d_2}e^{-\delta|u|/\ell_{\sigma_{\min}}}$, $d_1+d_2=-4$ (the product of two slices decays at the finest line); $f=\partial^aC^{(j_f)}(y-w)$, $a\in\{0,1,2\}$ ($a\le1$ for a compound as formed, \Cref{rot:h1s:nd}(ii); $a=2$ when the compound is the differenced factor of (R2b), see (iv)), $\sigma_{\max}\le j_f$; $K^F_\sigma$ the own-frame second-order remainder sums $K^X_\sigma(x,w)\defeq\sum_{y}\kappa_\sigma(x-y)T^X_xf^X(y)$, $K^Y_\sigma(\tau x,\tau w)\defeq\sum_{y'\in\Lambda_Y}\kappa_\sigma(R^{-1}(\tau x-y'))T^Y_{\tau x}f^Y(y')$, with lattice differences in the Taylor polynomials. Put $s\defeq\alpha/(1+\alpha)$, $\kappa\defeq\alpha/(2(1+\alpha))$.
\begin{enumerate}[label=\textup{(\roman*)},nosep,leftmargin=*]
\item $|\sum_{\sigma:\sigma_{\max}\le j_f}K_\sigma(x,w)|\le C\ell_{j_f}^{-2-a}e^{-\delta|x-w|/(4\ell_{j_f})}$ per flow: the compound is a line of dimension $-2-a$ at scale $j_f$. With an external or coarse leg $E$: $|\sum_{\sigma_{\max}\le j}K_\sigma(x)|\le C\ell_j^\alpha\|E\|$, a vertex weight $\alpha$ at scale $\sigma_{\max}$, summed geometrically.
\item $\sum_{\sigma_{\max}\le j_f}|K^X_\sigma(x,w)-K^Y_\sigma(\tau x,\tau w)|\le C(1+r)\eta'^{\theta_U}\ell_{j_f}^{-2-a-\theta_U}e^{-\delta|x-w|/(4\ell_{j_f})}$ with $\theta_U\defeq\min\{\theta,s/2^{a+1}\}$; with an external or coarse leg $\le C(1+r)\eta'^{\theta_U}\|E\|$ with $\theta_U\defeq\min\{\theta,\kappa s\}$. No frame vector appears.
\item For the double compound $(U_1,l,U_2)$ with $l=\partial^aC^{(j_l)}$, $a\in\{0,1\}$: per flow a line of dimension $-3-a$ at scale $j_l$; the difference gains $(\eta'/\ell_{j_l})^{s/8}$, with one logarithm per unit. The same holds when $l$ is a coarse element $\Gamma^F$: per flow $|\sum_{\sigma_1,\sigma_2}(U_1,\Gamma,U_2)(x,w)|\le Ce^{-m_r|x-w|/2}$ (dimension $0$ at the top scale), and the difference is $\le C(1+r)^2\eta^{\kappa s}e^{-m_r|x-w|/4}$.
\item First-order forms (the compounds $(T,f)$, $(SC,f)$ and the brackets of (R2b) whose differenced factor is a line carrying $a\le2$ derivatives): gain $(\eta'/\ell_{j_f})^{1/2^{a+2}}\ge(\eta'/\ell_{j_f})^{1/16}$ (slice leg) resp.\ $\eta^{1/4}$ (external leg), with one logarithm. When the differenced factor of (R2b) is a compound $(U,f)$, the bracket is bounded by (i)--(ii) with $a+1\le2$ in place of $a$: per flow $|t-x|$ times a line of dimension $-3-a$ at scale $j_f$, and the gain $\eta'^{\theta_U}$ with $\theta_U\ge\min\{\theta,s/8\}$.
\end{enumerate}
\end{lemma}
\begin{proof}
(i) $\sum_u\eta'^3|\kappa_\sigma(u)||u|^{1+\alpha}\le C\ell_{j_1}^{d_1}\ell_{j_2}^{d_2}\ell_{\sigma_{\min}}^{4+\alpha}\le C\ell_{\sigma_{\max}}^\alpha$ (the sum is dominated by $|u|\sim\ell_{\sigma_{\min}}$, where the coarser line is bounded by $\ell_{\sigma_{\max}}^{d_{\max}}$; more precisely it is $\le C\ell_{\sigma_{\min}}^{\alpha}(\ell_{\sigma_{\min}}/\ell_{\sigma_{\max}})^{|d_{\max}|}$, geometric in $\sigma_{\min}$ at fixed $\sigma_{\max}$); $|T_xf(y)|\le C\ell_{j_f}^{-2-a-\alpha}|u|^{1+\alpha}$ with the leg's decay $e^{-\delta|y-w|/(2\ell_{j_f})}\le e^{-\delta|x-w|/(4\ell_{j_f})}e^{\delta|u|/(4\ell_{j_f})}$, the last factor absorbed by the unit's decay at $\ell_{\sigma_{\min}}\le\ell_{j_f}$; $\sum_{\sigma_{\max}\le j_f}\sum_{\sigma_{\min}\le\sigma_{\max}}\ell_{\sigma_{\min}}^{\alpha}(\ell_{\sigma_{\min}}/\ell_{\sigma_{\max}})^{|d_{\max}|}\le C\ell_{j_f}^\alpha$. External leg: the same with $\|E\|$. (ii) Split $K^X_\sigma-K^Y_\sigma=\sum_y[\kappa_\sigma(x-y)-\kappa_\sigma(R^{-1}(\tau x-\tau y))]T^X_xf^X(y)+\sum_y\kappa_\sigma(R^{-1}(\tau x-\tau y))[T^X_xf^X(y)-T^Y_{\tau x}f^Y(\tau y)]=:A_\sigma+B_\sigma$. $A_\sigma$: telescoping the product of the unit's lines with \Cref{rot:h1s:P} ($n\le2$) for the differenced line and \eqref{eq:h1s:B} for the other gives $|\Delta\kappa_\sigma|\le C(\eta'/\ell_{\sigma_{\min}})^{1/4}\ell_{j_1}^{d_1}\ell_{j_2}^{d_2}e^{-\delta|u|/(2\ell_{\sigma_{\min}})}$ (with the $125^\theta$-form for $\sigma_{\min}\le2$), so as in (i) $\sum_{\sigma_{\max}\le j_f}A_\sigma\le C\eta'^\theta\ell_{j_f}^{\alpha-\theta}\ell_{j_f}^{-2-a-\alpha}$ for $\theta\le\frac14$, $\theta<\alpha$. $B_\sigma$: the bracket is $\le\min\{P,Q\}$ with $P\defeq C\ell_{j_f}^{-2-a-\alpha}(|u|^{1+\alpha}+\eta'^{1+\alpha})$ and, by \Cref{rot:h1s:E} for a slice leg, $Q\defeq b_0+b_1|u|+b_2$ with $b_0=C(\eta'/\ell_f)^{1/2^a}\ell_f^{-1-a}$, $b_1=C(\eta'/\ell_f)^{1/2^{a+1}}\ell_f^{-2-a}$, $b_2=C\eta'\ell_f^{-2-a}$ ($\ell_f\defeq\ell_{j_f}$). Since $(1+\alpha)(1-s)=1$, $\min\{P,Q\}\le P^{1-s}Q^s\le C\gamma_f\ell_f^{-2-a}(|u|+\eta')(1+|u|/\ell_f)^s$ with $\gamma_f\defeq(\eta'/\ell_f)^{s/2^{a+1}}$ (the $\ell_f$-exponent is $-(2+a+\alpha)(1-s)-(1+a)s=-2-a$); then $B_\sigma\le C\gamma_f\ell_f^{-2-a}e^{\dots}\sum_u\eta'^3|\kappa_\sigma(u)|(|u|+\eta')$, and the first-order sum of a degree-$(-1)$ kernel is $\le C$ per $\sigma$, hence $\le C(1+r)$ after the $\sigma$-sum: one logarithm. External leg: $b_0=C\eta^{1/2}\|E\|$, $b_1=C\eta^\kappa\|E\|$, $b_2=C\eta'\|E\|$, $\gamma_f=\eta^{\kappa s}\le5\eta'^{\kappa s}$. (iii) Slice $l$: $|T_xT_wl(y,z)|\le\ell_l^{-3-a-2\alpha}|u_1|^{1+\alpha}|u_2|^{1+\alpha}$ (mixed second-order lattice remainder, from \eqref{eq:h1s:B} at order $a+4$) gives, summed, $\ell_l^{-3-a}$ per flow. Difference: $P=\ell_l^{-3-a-2\alpha}|u_1|^{1+\alpha}|u_2|^{1+\alpha}$ and $Q\le C(\eta'/\ell_l)^{1/8}\ell_l^{-1-a}(1+|u_1|/\ell_l)(1+|u_2|/\ell_l)$, since the differenced factors $l,\nabla l,\nabla\nabla l$ carry $n\le a+2\le3$ derivatives (\Cref{rot:h1s:P}, exponent $\frac18$) plus the displacements; $P^{1-s}Q^s\le C(\eta'/\ell_l)^{s/8}\ell_l^{-3-a}|u_1||u_2|(\dots)^s$ as $(3+a+2\alpha)(1-s)+(1+a)s=3+a$; two first-order sums give $(1+r)^2$. Coarse $l=\Gamma^F=H^FB^FH^{F*}$ (\Cref{rot:lem:Gam}): $T_xT_w$ acts on the two $H$-factors separately, each $C^{1,\alpha}$ in its first variable with $|T_xH^X(\cdot,p)(y)|\le C|y-x|^{1+\alpha}e^{-\delta|x-p|}$, and $|B^F(p,q)|\le Ce^{-\delta|p-q|}$, so $|T_xT_w\Gamma^F(y,z)|\le C|u_1|^{1+\alpha}|u_2|^{1+\alpha}e^{-\delta|x-w|/2}$; the difference uses $|\Gamma^X-\Gamma^Y\circ\tau|\le C\eta^{1/2}$ and \Cref{rot:h1s:E} for each $H$-factor, $Q\le C\eta^\kappa(1+|u_1|)(1+|u_2|)$, whence $P^{1-s}Q^s\le C\eta^{\kappa s}|u_1||u_2|$. (iv) First-order forms: $P=\ell_f^{-2-a}|u|$ (slice; $C\|E\||u|$ external), $Q$ as above with $n\le a+1$ in the gradient; $P^{1/2}Q^{1/2}\le C(\eta'/\ell_f)^{1/2^{a+2}}\ell_f^{-1-a}|u|^{1/2}(\dots)$ resp. $C\eta^{1/4}\|E\||u|^{1/2}(\dots)$, and the first-order sum converges per $\sigma$ with the $(1+r)$ after the $\sigma$-sum. For the brackets of (R2b) the switched factor $\Delta[g(t)-g(x)]$ is $\le|t-x|$ times the slice law for $\nabla g$ plus $2c_\tau\eta'\sup|\nabla g^Y|$: the degree-one term with its gain and the explicit-$\eta'$ term (degree zero, one logarithm). When the differenced factor is a compound $(U,f)$ at an outer vertex $t\in\{x,w\}$, its lattice difference is the compound with the differenced leg: at $t=w$ directly, and at the unit vertex $t=x$ by the translation invariance of $\kappa_\sigma$, since shifting the summation variable gives $K_\sigma(x+\eta'e_\mu,w)-K_\sigma(x,w)=\eta'\sum_y\kappa_\sigma(x-y)\,T_x(\partial_\mu f)(y)$ exactly, the lattice differences commuting (and likewise on $\Lambda_Y$ with the step $\eta'Re_\mu$). So the difference of the bracket is a sum of $|t-x|/\eta'$ terms $\eta'\,[K^X_\sigma-K^Y_\sigma]$ with the leg $\partial_\mu f$, $a+1\le2$, each bounded by (ii) (\Cref{rot:h1s:E} for the slice leg with $a+1$, \Cref{rot:h1s:P} at $n\le3$; the unit lines keep $n\le2$): per flow $|t-x|\,\ell_{j_f}^{-3-a}$, gain $\eta'^{\theta_U}$ with $\theta_U=\min\{\theta,s/2^{a+2}\}\ge\min\{\theta,s/8\}$, one logarithm; the $Y$-displacement term is explicit-$\eta'$ as before.
\end{proof}

% r03_h1s_e: proof of the two-flow comparison bound, part 5: the telescoping, every term is small, the exponent theta_1, the sectors, the proof of the theorem.

\paragraph{The telescoping.}
After (NF1)--(NF4) and (R1)--(R3) in each flow's own frame, $E^F(\cR)$ for one ordering class is a finite sum of multilinear expressions in the elements: slice lines and top pieces (with at most two lattice derivatives, \Cref{rot:h1s:nd}), compounds $(U,f)$ and $(U_1,l,U_2)$ (own-frame remainder sums, \Cref{rot:h1s:U}), first-order compounds $(T,f)$, $(SC,f)$, the brackets $[G(t)-G(x)]$ of (R2b) with one differenced factor, the coarse elements $\Gamma^F$, $H^F$, $B^F$, the externals and vertex functions at $\tau$-matched points (shifts included), the contour elements at $e>0$, and the common numbers (the zero moments $c$, the order-four shares, the vacuum-polarisation leftover weight, the one-point structures), which are never switched. The two flows' expressions have the same combinatorics. Order the elements consistently (compounds first, then lines, coarse elements, externals, contour elements) and telescope:
\[
E^X(\cR)-E^Y(\cR)=\sum_iT_i,\qquad T_i\defeq\big[e_1,\dots,e_{i-1}\text{ in the $Y\circ\tau$ version};\ e_i\to e_i^X-e_i^Y\circ\tau;\ e_{i+1},\dots\text{ in the $X$ version}\big],
\]
an identity, every term being the same multilinear form with elements substituted, the co-switched pairs (a self-energy and its counterterm share, the two Wick contractions of \Cref{rot:h1s:M}, a scalar W and its order-four share) switched together; at most $n_\cR$ terms per class, and finitely many classes per family. Compounds are switched first so that a unit and its leg are compared as a whole and no mixed-frame first moment is ever written; externals are switched late so that the only frame-sensitive quantity left is the difference of the two flows' own-frame remainders, which \Cref{rot:h1s:E} bounds. Identities (the counterterm cancellations, the Ward sum, the moment subtractions) are used only at the two pure endpoints, in each flow's own frame; the mixed terms are bounded in absolute value.

\begin{lemma}[every telescoping term is small]\label{rot:h1s:T}
For every ordering class and every term $T_i$, with $\theta\le\theta_1=\alpha^2/(4(1+\alpha)^2)$,
\[
|T_i|\le C(1+k)^p\eta'^\theta\cdot[\text{the per-flow bound \cite[(1.33)]{Ba3} of the monomial}],\qquad p=p(\bar n),
\]
with the prefactor $p_r^2\delta_r^{-1/2}$ when $e_i$ is a contour element.
\end{lemma}
\begin{proof}
Apply \Cref{rot:h1s:MS} to $|T_i|$, the absolute value of a multilinear form whose degrees are those of \Cref{rot:h1s:Dp}(a).
(i) \emph{Per-flow bounds} of every element in $X$-coordinates, whichever flow it comes from: lines by \eqref{eq:h1s:B}, with $\delta/2$ for the $Y\circ\tau$ versions; compounds by \Cref{rot:h1s:U}(i),(iii), lines of dimension $-2-a$, resp.\ $-3-a$ ($0$ at the top scale for a $\Gamma$ middle), at their leg's scale, or vertex weights $\alpha$ for external legs; first-order compounds by \Cref{rot:h1s:U}(iv) (degree-one objects); the brackets of (R2b) by $|G(t)-G(x)|\le|t-x|\sup|\nabla g|$ plus the explicit-$\eta'$ displacement term; coarse elements, externals and vertex functions by their table; the common numbers by \Cref{rot:h1s:Z0} and \Cref{rot:h1s:MS} per flow ($\le C(1+r)^{n_0+1}$ each; the vacuum-polarisation leftover $\le C\ell_J^{-1}$ as a weight $-1$); at $e>0$ the contour elements by \Cref{rot:prop:BD}(b), a contour line that received a derivative being a line of dimension $-1$ at its slice (the contour sum carries the weight $\eta$ per bond against $\sim\ell_j/\eta$ bonds within the decay range, \cite[(3.63)--(3.64)]{K1} with the contour clause of \cite[Prop.~3.7]{K1}). Decay is re-anchored at merged vertices as in the first step of \Cref{rot:h1s:MS}.
(ii) \emph{The switched element}, divided by $C\eta'^\theta$, obeys the bound of its type weakened by $\ell_j^{-\theta}$ at its scale (or unweakened if it is fixed): a line with $n\le2$ derivatives by \Cref{rot:h1s:P}, gain $(\eta'/\ell_j)^{1/4}\ge(\eta'/\ell_j)^\theta$ (top piece: $\eta'^{1/4}$); a compound by \Cref{rot:h1s:U}(ii),(iii), gain $\eta'^{\theta_U}\ell_{j_f}^{-\theta_U}$ with $\theta_U\ge\theta$; a first-order compound or a bracket by \Cref{rot:h1s:U}(iv), gain $\ge(\eta'/\ell_j)^{1/16}$ resp.\ $\eta^{1/4}$ (a bracket whose differenced factor is a compound: $\eta'^{\theta_U}\ell_{j_f}^{-\theta_U}$ with $\theta_U\ge\min\{\theta,s/8\}$); a coarse element by $C\eta^{1/2}$; an external by $C\eta^{1/2}\|E\|$, its gradient as a recipient by $C\eta^\kappa\|E\|$ (\Cref{rot:h1s:E}); a vertex function by $c_\tau\eta'\|\nabla w\|$, and the $\partial^{\eta'}g$ term of (R2a), of size $O(\eta')$ in both flows, by $O(\eta')\le\eta'^\theta\ell_{j(v)}^{1-\theta}$; at $e>0$ a contour element by \Cref{rot:prop:BD}(b): $(\eta'/\ell_j)^{1/2}$ for contour lines, also with the one lattice derivative a recipient contour line carries, $\eta^{1/2}$ for contour vertices, $e\,p_r\lmu^{-1}c_\tau\delta_r^{-1/2}\eta'$ for the phases, whence the prefactor; contour vertices are never recipients (\Cref{rot:h1s:rule}(iii)). For $j\le2$ there is no gain and the factor is paid by the degree loss $\theta$, as in \Cref{rot:cor:unif}(a).
(iii) \emph{Degrees.} By \Cref{rot:h1s:Dp} every component of every initial segment of the reduced ordering has $D\ge1$, or $D\in\alpha+\N_0$, or is an explicit-$\eta'$ component with $D=1-\theta$ after $\eta'^\theta$ is taken out; fixed elements impose no condition. Hence $D(H)-\theta>0$ for every component containing the switched element, and the scale sums of \Cref{rot:h1s:MS} converge geometrically there.
(iv) \emph{Logarithms.} $n_0\le\bar n$ zero-degree components, $(1+r)$ per compound, $(1+r)^{n_0+1}$ per common number, $(1+r)^2$ per double compound, and the $(1+r)^{n_0}$ of \Cref{rot:h1s:MS} itself; with $\le\bar n$ elements, $p(\bar n)\le3\bar n+2$ suffices, and $(1+r)\le(1+k)$.
\Cref{rot:h1s:MS} then gives $|T_i|\le C(1+k)^p\prod_lc_l$ with $\prod_lc_l=\eta'^\theta\cdot[\text{per-flow constants}]$; the localisation \cite[(3.55)]{K1} costs at most $p_r^3$ per coarse line, absorbed in the field norms of step (5) of the proof of \Cref{rot:thm:A}.
\end{proof}

\paragraph{The exponent $\theta_1$.}
In every term the switched element carries $\eta'^{\theta'}\ell_{j}^{-\theta'}$ with its own exponent $\theta'$, and \Cref{rot:h1s:MS} at the stage where the component $H\ni e_i$ is shrunk sums $\sum_{j\le j'}\ell_j^{D(H)-\theta}\eta'^\theta\le C\ell_{j'}^{D(H)-\theta}\eta'^\theta$, geometric because $D(H)-\theta>0$; the remaining factor is the merged vertex's weight, and since all later components containing it again have $D-\theta>0$ the factor $\eta'^\theta$ survives. (When $D(H)=0$ the sum $\sum_j(\eta'/\ell_j)^\theta=\sum_j5^{-j\theta}$ is $O(1)$ and the gain is lost; \Cref{rot:h1s:Dp} excludes this.) With $s=\alpha/(1+\alpha)$, $\kappa=\alpha/(2(1+\alpha))$, the gains are: a switched line with $n\le2$ derivatives, $1/2^n\ge\frac14$; the top piece, $\frac14$; a compound with a slice leg carrying $a\le2$ derivatives ($a=2$ only for the differenced compound of (R2b)), $s/2^{a+1}\ge s/8$; with an external or coarse leg, $\kappa s$; the double compound, $s/8$ (slice middle) or $\kappa s$ (coarse middle); the first-order forms, $1/2^{a+2}\ge\frac1{16}$ (slice) or $\frac14$ (external); coarse elements, $\frac12$; vertex functions and explicit-$\eta'$ terms, $1$; contour lines, $\frac12$; contour vertices, $\frac12$; phases, $1$ with the prefactor. The arithmetic: $\kappa s=\alpha^2/(2(1+\alpha)^2)=2\theta_1$; $s/4-\theta_1=\alpha/(4(1+\alpha)^2)>0$; $s/8-\theta_1=\alpha(1-\alpha)/(8(1+\alpha)^2)>0$ for $\alpha<1$; $\frac1{16}-\theta_1=(1-\alpha)(1+3\alpha)/(16(1+\alpha)^2)>0$; and $\min(\alpha,\frac12)-\theta_1>0$ because $\theta_1<\alpha/4\le\alpha$ and $\theta_1<\frac1{16}$. So $\theta_1\defeq\alpha^2/(4(1+\alpha)^2)$ is at most every gain exponent (the weakest two being $s/8$ and $2\theta_1$), and every component containing the switched element keeps the margin $D-\theta_1\ge\min(\alpha,1)-\theta_1>0$. Nothing is nested: no compound is the leg of another (the double compound is one element, and the recipients of (R2) receive a single gradient), so $\theta_1$ does not depend on $\bar n$. The consumers (\Cref{rot:lem:r}, \Cref{rot:prop:cp}) use only $\theta_1>0$ and a polynomial $(1+k)^p$, so the exact value plays no role; what matters is that the set of element types and derivative counts is finite, that every type has a difference law with some positive exponent, and that the local terms are common numbers.

\begin{proof}[Proof of \Cref{rot:thm:H1S}]
Sum \Cref{rot:h1s:T} over the $\le n_\cR$ terms of each ordering class and over the finitely many classes (\cite[(2.7)]{Ba3}; their number depends on the family only). The sectors:
\emph{$e>0$.} The constant-background reduction (\Cref{rot:lem:CB}) is flow-independent algebra; after it the new elements are the contour elements of \Cref{rot:prop:BD}(b). Every connected subgraph containing a contour vertex has degree $\ge\frac12$ (\cite[(3.30)]{K1}; \Cref{rot:lem:CB}(3)), so no new structure of degree $\le0$ appears and \Cref{rot:h1s:Dp} is unchanged; a contour line may be the recipient of (R2a)(ii) at its bond end, and this is the one-derivative clause of \Cref{rot:prop:BD}(b); the phases are Lipschitz in the vertex position with the bound $e\,p_r\lmu^{-1}\delta_r^{-1/2}$, inside the prefactor. $e$ enters only through $e_r^2=e^2\delta_r$, so the bound holds at every $e>0$.
\emph{The sectors with legs at $A=0$.} Scalar two-point: orders $\lambda$, $\lambda^2$, $e^2$, $e^4$ with the counterterm cancellations of \Cref{rot:h1s:VP}(d), $e^2\lambda$ trivial; the $(2,1)$-sector at order $e^3$: T and SC as first-order compounds (no Ward identity of the vertex class is used: in the per-flow normal form the identities of this sector are identities on all-$C$ expressions, applied in each flow alone, and a reduction of this sector to a ``mixed zero-moment'' bound for current--current subgraphs with lines from different flows is not used, the mixed zero moment being a lattice-scale quantity of order $e_r^2/\eta'$); the vacuum polarisation at $e^2$ through \Cref{rot:h1s:VP}(a),(b) and at $e^4$ through (c); L24; everything else of positive degree (\cite[Cor.~2.3, (2.17)]{Ba3}).
\emph{The vacuum sector.} All-$C$ vacuum monomials cancel against $E_1$ exactly (\Cref{rot:lem:ctsame}, \cite[(3.35)]{K1}); the $\Gamma$-monomials, cut at their $\Gamma$-lines, leave all-$C$ components with at least two legs, renormalised by the rule with the $\Gamma$-lines as legs and recipients (one-point components are common numbers, (R2c)); \Cref{rot:h1s:T} applies per coarse site, times $|T|\delta_r^{-3}$ sites (\Cref{rot:thm:V}(e)).
\emph{The source vertex.} The $h$-vertex has degree $2$ (\Cref{rot:lem:S1}); the $h$-families other than $\cT_h$ contain no new structure of degree $\le0$.
Hence $|E^X(\cR)-E^Y(\cR)|\le Cn_\cR\,p_r^2\delta_r^{-1/2}\eta'^\theta(1+k)^p\cdot[\text{right side of \cite[(1.33)]{Ba3}}]$ for $0<\theta\le\theta_1$, uniformly in $r$, $k$, the torus (the tori enter through periodisation, \Cref{rot:h1s:Z0}) and the masses, and in every term of the telescoping the switched element obeys the bound of its type weakened by $\theta$, which is the statement.
\end{proof}

